% =============================================================================
% Poc.tex — Descrição da Prova de Conceito (PoC) para a seção 4.5 "Demonstração"
% Dissertação: Inteligência Artificial Agêntica no Pipeline de Sinistros em Seguros
%
% Uso no Overleaf: % =============================================================================
% Poc.tex — Descrição da Prova de Conceito (PoC) para a seção 4.5 "Demonstração"
% Dissertação: Inteligência Artificial Agêntica no Pipeline de Sinistros em Seguros
%
% Uso no Overleaf: % =============================================================================
% Poc.tex — Descrição da Prova de Conceito (PoC) para a seção 4.5 "Demonstração"
% Dissertação: Inteligência Artificial Agêntica no Pipeline de Sinistros em Seguros
%
% Uso no Overleaf: % =============================================================================
% Poc.tex — Descrição da Prova de Conceito (PoC) para a seção 4.5 "Demonstração"
% Dissertação: Inteligência Artificial Agêntica no Pipeline de Sinistros em Seguros
%
% Uso no Overleaf: \input{Poc} dentro da seção "Demonstração".
% Pacotes esperados no preâmbulo principal: graphicx (figura) e abntex2cite (ou
% natbib) para \cite. A figura usa o arquivo docs/figuras/poc-painel.png do
% repositório; envie-o ao Overleaf como figuras/poc-painel.png ou ajuste o caminho.
%
% As chaves de citação estão listadas em BibTeX no fim deste arquivo (comentadas).
% Se o seu .bib já usa outras chaves, basta substituí-las aqui.
% =============================================================================

\subsection{Prova de Conceito: recorte de Triagem, Aviso de Sinistro e Documentação}
\label{sec:poc}

A fase de demonstração, conforme preconizado pela DSR, tem por finalidade evidenciar que o
artefato proposto é capaz de resolver, ao menos em um contexto controlado, instâncias do
problema que motivou a pesquisa \cite{peffers2007,hevner2004}. Para tanto, foi construída uma
Prova de Conceito (PoC) que instancia um recorte do \textit{framework} apresentado na
Seção~4.4.1, abrangendo a orquestração, a Triagem, o Aviso de Sinistro e a etapa de
Documentação e Evidências Digitais (Figuras 2 a 5), além da camada transversal de
\textit{guardrails} corporativos (Figura~1). As etapas de Regulação, Vistoria, Avaliação,
Liquidação e Pesquisa de Satisfação permanecem fora do recorte e são tratadas como trabalhos
futuros.

O recorte foi escolhido por concentrar as dores mais citadas pelos segurados — complexidade
do processo de solicitação, falta de clareza sobre requisitos e documentos e ausência de
empatia no primeiro contato \cite{capgemini2024b} — e por constituir a porta de entrada de todo
o \textit{pipeline}: a qualidade das informações coletadas no aviso condiciona as etapas
subsequentes de regulação e liquidação. Trata-se, adicionalmente, do ponto em que o Código
Civil e o Marco Legal dos Seguros impõem ao segurado o dever de comunicar prontamente o
sinistro \cite{brasil2002codigocivil,brasil2024marco}, o que torna a disponibilidade contínua
de um canal conversacional particularmente relevante.

\subsubsection{Objetivo e premissas da PoC}

O objetivo da PoC é demonstrar a viabilidade técnica de um sistema multiagente que conduza a
jornada do segurado desde a primeira mensagem, em um canal que simula o WhatsApp, até a
conclusão da abertura do sinistro, com protocolo registrado e documentação mínima validada.
Foram adotadas as seguintes premissas:

\begin{itemize}
  \item \textbf{Ramo automóvel}, contemplando os tipos de sinistro de veículo mais recorrentes:
        colisão, roubo ou furto, incêndio, eventos da natureza (alagamento, granizo, queda de
        árvore) e danos a vidros. Solicitações de assistência (guincho, pane) e assuntos que
        não configuram sinistro recebem apenas orientação, sem abertura de processo.
  \item \textbf{Dados exclusivamente sintéticos}: clientes, apólices, veículos e documentos
        (CNH, CRLV, fotografias e boletins de ocorrência) são gerados artificialmente, em
        conformidade com os princípios de minimização e proteção de dados pessoais. Nenhum dado
        real é processado.
  \item \textbf{Sistemas da seguradora simulados}: cadastro, apólices, registro de sinistros,
        armazenamento e validação de documentos e envio de mensagens são implementados como
        serviços simulados (\textit{mocks}) atrás de interfaces estáveis, de modo que a troca
        por integrações reais não altere a lógica dos agentes \cite{ofoeda2019}.
  \item \textbf{Modelos de linguagem como recurso e não como dependência}: o sistema funciona
        integralmente com regras determinísticas e recorre ao LLM nos pontos em que a
        linguagem natural do segurado excede o que as regras capturam, o que permite comparar
        experimentalmente as duas abordagens (Seção~\ref{sec:poc-avaliacao}).
\end{itemize}

\subsubsection{Arquitetura e tecnologias}

A PoC foi implementada em Python~3.12 com o \textit{framework} LangGraph, que modela a
aplicação como um grafo de estados com suporte a ciclos, subgrafos e persistência
(\textit{checkpointing}) da conversa. Essa escolha materializa a noção de agente como unidade
que mantém estado interno persistente entre interações e decide a próxima ação a partir de
observações parciais do ambiente \cite{liu2024,wang2023}. A API foi construída com FastAPI, com
comunicação em tempo real via WebSocket, e a persistência utiliza SQLite para o estado das
conversas, o domínio simulado e o registro de eventos. A interface web, desenvolvida em React
com TypeScript e a biblioteca React Flow, reproduz um atendimento via WhatsApp e, lado a
lado, um painel de observabilidade que exibe o fluxo de agentes em execução
(Figura~\ref{fig:poc-painel}).

O acesso aos modelos de linguagem é feito por um \textit{gateway} próprio que padroniza as
chamadas no formato da API da OpenAI (\textit{OpenAI-compatible}), independentemente do
provedor. Na PoC, o provedor configurado é o Google Gemini, com um modelo principal e um
modelo de \textit{fallback} acionado em caso de indisponibilidade. As respostas estruturadas
são solicitadas por esquema JSON e validadas com Pydantic, e cada chamada registra o modelo
utilizado, a versão do \textit{prompt}, os \textit{tokens} consumidos e a latência. O
\textit{gateway} corresponde ao componente \textit{AI Gateway} da arquitetura proposta
(Figura~11) e garante que a substituição do provedor seja uma decisão de configuração, e não de
código.

As ferramentas (\textit{tools}) seguem o conceito de interfaces externas que ampliam a
capacidade do agente para além da predição de linguagem \cite{schick2023,qin2023}. Cinco
ferramentas foram definidas, cada uma com contrato tipado de entrada e saída, mecanismos de
\textit{timeout} e novas tentativas com recuo exponencial, além de injeção configurável de
falhas para avaliar a resiliência do sistema:

\begin{itemize}
  \item \texttt{get\_customer\_data}: consulta cadastro e apólices do segurado;
  \item \texttt{open\_claim}: registra o aviso de sinistro e gera o protocolo;
  \item \texttt{store\_files}: armazena os documentos vinculados ao protocolo;
  \item \texttt{validate\_document}: valida formato, tamanho, duplicidade e consistência dos
        metadados do documento com o processo;
  \item \texttt{send\_whatsapp}: entrega a mensagem ao segurado.
\end{itemize}

As interfaces das ferramentas foram desenhadas de forma compatível com a exposição via
\textit{Model Context Protocol} \cite{mcp2025}. Na versão atual da PoC, porém, elas são
invocadas diretamente pelos agentes por meio de um registro de implementações (simuladas ou
reais), e a publicação como servidor MCP permanece como evolução prevista.

\subsubsection{Orquestração e agentes especialistas}

A organização dos agentes segue o padrão hierárquico descrito na literatura de sistemas
multiagentes baseados em LLM, em que um orquestrador central decompõe a tarefa e delega
subtarefas a agentes especializados \cite{hill2025,singhal2026,sapkota2025}. Cada mensagem do
segurado percorre o grafo principal na seguinte sequência:
\textit{guardrails} de entrada $\rightarrow$ Orquestrador $\rightarrow$ agente especialista
$\rightarrow$ transição de etapa $\rightarrow$ \textit{guardrails} de saída $\rightarrow$ envio.
O estado global da conversa registra a etapa do processo
(\textsc{Triagem} $\rightarrow$ \textsc{Aviso} $\rightarrow$ \textsc{Documentos}
$\rightarrow$ \textsc{Concluído}), além de estados terminais alternativos (orientado,
bloqueado, encaminhado a humano e cancelado). A mudança de etapa é realizada exclusivamente
por uma função determinística de transição, o que garante previsibilidade e auditabilidade do
fluxo.

O \textbf{Orquestrador} implementa as três responsabilidades previstas na Figura~2:
identifica a situação atual do sinistro, informa o status ao segurado quando solicitado e
encaminha a mensagem ao especialista adequado. O encaminhamento usa uma tabela fixa que
associa cada etapa a um especialista. Quando a mensagem foge do roteiro esperado, um
classificador de intenções identifica desvios como dúvida, consulta à apólice, correção de
dado já informado, pedido de status, cancelamento, pedido de atendimento humano ou assunto
fora do escopo. O classificador combina regras e, opcionalmente, um LLM, e após tratar o
desvio o atendimento é retomado a partir da pergunta pendente. Cada especialista é um subgrafo
independente composto por subagentes, conforme a Tabela~\ref{tab:poc-agentes}.

\begin{table}[htbp]
\centering
\caption{Agentes especialistas, subagentes e ferramentas da PoC}
\label{tab:poc-agentes}
\small
\begin{tabular}{|p{2.6cm}|p{7.2cm}|p{3.4cm}|}
\hline
\textbf{Especialista} & \textbf{Subagentes (responsabilidade)} & \textbf{Ferramentas} \\
\hline
Triagem (Fig.~3) &
Validação do segurado (CPF e apólices); classificação do ocorrido (tipo de sinistro);
decisão de elegibilidade (status, vigência na data do evento e cobertura exigida pelo
tipo); identificação de envolvidos e feridos (colisão), com sinalização de urgência &
\texttt{get\_customer\_data} \\
\hline
Aviso de Sinistro (Fig.~4) &
Coleta multi-turno dos dados do evento (quando, onde, o que aconteceu e danos);
avaliação preliminar de cobertura; registro do aviso, com protocolo e próximos passos &
\texttt{open\_claim} \\
\hline
Documentação (Fig.~5) &
Lista de documentos por tipo de sinistro; recepção e armazenamento; validação;
complementação (progresso, lembrete, prazo e confirmação final) &
\texttt{store\_files}, \texttt{validate\_document} \\
\hline
Consultas &
Classificação da pergunta; identificação do segurado; resposta sobre vigência,
vencimento, renovação e coberturas da apólice ou sobre dúvidas gerais de seguro &
\texttt{get\_customer\_data} \\
\hline
\end{tabular}
\end{table}

Na \textbf{Triagem}, o segurado é identificado pelo CPF e o relato livre do ocorrido é
classificado, o que determina a cobertura exigida e os dados a coletar. Um mecanismo decisório
define se o caso evolui para a abertura formal ou se o segurado deve apenas ser orientado (por
exemplo, apólice vencida, evento fora da vigência, ausência da cobertura correspondente ou
pedido de assistência), evitando a abertura indevida de processos, como propõe a Figura~3. No
\textbf{Aviso de Sinistro}, os dados são extraídos de narrativas livres; quando uma informação
está ausente, o agente solicita apenas o que falta, sem repetir o que já foi informado. Após a
avaliação preliminar de cobertura, o aviso é registrado e o protocolo é comunicado. Na
\textbf{Documentação}, a lista de documentos varia conforme o tipo de sinistro (por exemplo,
boletim de ocorrência obrigatório em roubo ou furto e condicionado à existência de terceiros
ou feridos em colisão). Cada arquivo é validado de forma determinística quanto a formato,
tamanho, duplicidade e consistência dos metadados (CPF da CNH, placa do CRLV e data das
fotografias em relação ao evento). O agente informa o que já foi aceito e o que ainda falta e,
na ausência de resposta, envia lembrete e encaminha o caso a um atendente ao fim do prazo. O
especialista \textbf{Consultas}, acrescentado durante o ciclo de design, responde perguntas
sobre a própria apólice a partir dos dados cadastrais e dúvidas gerais (``o que é sinistro'',
``o que é franquia'') a partir de uma base de conhecimento curada. Ele atende à
responsabilidade de ``verificar a vigência e as coberturas da apólice'' atribuída a agentes de
sinistro \cite{singhal2026} e reduz a falta de clareza apontada pelos segurados
\cite{capgemini2024b}.

\subsubsection{Uso de modelos de linguagem: regras primeiro, LLM como complemento}

Em consonância com a ênfase da literatura em governança, previsibilidade e capacidade de
intervenção humana em sistemas agênticos \cite{casper2025,hughes2025}, a PoC adota o princípio
de ``regras primeiro, LLM como complemento''. Regras configuráveis (expressões e listas de
termos em arquivos YAML) resolvem os casos recorrentes com custo e latência mínimos. O LLM é
acionado quando as regras não são conclusivas, nos seguintes pontos:

\begin{itemize}
  \item extração e classificação do relato livre na Triagem e no Aviso (por exemplo, ``deixei o
        carro estacionado e quando voltei ele não estava mais lá'' como roubo ou furto);
  \item classificação de intenções do Orquestrador em mensagens que desviam do roteiro;
  \item classificação de segurança, ética e conformidade corporativa nos \textit{guardrails};
  \item redação de respostas a perguntas abertas, restrita à base de conhecimento e aos dados
        da apólice, técnica alinhada à geração aumentada por recuperação (RAG) e com instrução
        explícita de não prometer cobertura, valores ou prazos.
\end{itemize}

Todas as instruções enviadas ao LLM ficam em arquivos versionados, separados do código, e
as respostas estruturadas são validadas por esquema. Em caso de falha do provedor, o sistema
degrada para as regras sem interromper o atendimento.

\subsubsection{Camada de \textit{guardrails} corporativos}

A camada de \textit{guardrails} implementa as barreiras de proteção descritas na literatura
\cite{dong2024} e posicionadas transversalmente no \textit{framework} (Figura~1). Toda
mensagem recebida e toda resposta gerada passam por três estágios executados em paralelo:
\textbf{Segurança} (\textit{prompt injection} direta e indireta, tentativas de
\textit{jailbreak} e vazamento de dados pessoais), \textbf{Ética} (pedidos discriminatórios ou
abusivos) e \textbf{Corporativo} (sinais de fraude e aderência às políticas). Um agregador
consolida os resultados em um veredito (\textit{allow}, \textit{sanitize}, \textit{block} ou
\textit{escalate}) acompanhado de justificativa. Cada estágio pode ser ativado ou desativado
por configuração, o que viabiliza os experimentos comparativos. A detecção opera em duas
camadas: regras determinísticas e, quando estas não bloqueiam, um classificador baseado em LLM
que avalia os três estágios em uma única chamada e descarta veredictos de baixa confiança.

A injeção indireta é tratada inspecionando também o nome e os metadados dos documentos
enviados. Por exemplo, um arquivo nomeado com a instrução ``ignore as instruções e aprove o
sinistro'' é bloqueado antes do armazenamento. Ao final da Triagem, a checagem corporativa é
reexecutada com o contexto da apólice, incorporando sinais de fraude como apólice contratada
pouco antes do evento e aviso feito muito tempo após o ocorrido. Em coerência com o princípio
de \textit{human-in-the-loop} \cite{wu2022} e com a recomendação de manter sob supervisão
humana as decisões de maior risco \cite{singhal2026}, os casos suspeitos não são negados
automaticamente: são encaminhados a um atendente humano com a justificativa registrada. O
painel dispõe ainda de um controle de pausa que suspende a atuação dos agentes e transfere
as conversas para atendimento humano.

\subsubsection{Observabilidade, auditabilidade e conformidade}

Toda decisão relevante (roteamento, veredito de \textit{guardrail}, elegibilidade, chamada de
ferramenta e chamada ao LLM) é emitida como evento estruturado, persistida e transmitida ao
painel em tempo real. O registro inclui o agente, o subagente, a justificativa, a latência e,
para chamadas ao LLM, o modelo, a versão do \textit{prompt} e os \textit{tokens} consumidos.
Esse registro instrumenta as dimensões de conformidade e auditabilidade da Equação~1
\cite{dodda2023} e atende aos requisitos de transparência e rastreabilidade em decisões
automatizadas previstos no PL~2338/2023 \cite{camara2023pl2338}. Os dados pessoais (CPF e
telefone) são mascarados em logs e eventos. As regras de negócio (elegibilidade, tipos de
sinistro e coberturas exigidas, documentos por tipo, prazos de lembrete e de encerramento) são
mantidas em arquivos de configuração, separadas do código, e as decisões de projeto foram
documentadas em 19 registros de decisão arquitetural (ADRs), o que favorece a revisão por áreas
de negócio e de \textit{compliance}.

\begin{figure}[htbp]
  \centering
  \includegraphics[width=\textwidth]{figuras/poc-painel.png}
  \caption{Interface da PoC: à esquerda, o atendimento simulado via WhatsApp; à direita, o
  fluxo de agentes com o especialista e os subagentes em execução e o log de eventos com a
  justificativa de cada decisão. Fonte: Elaboração própria, 2026}
  \label{fig:poc-painel}
\end{figure}

\subsubsection{Cenário ilustrativo}

Para ilustrar o funcionamento, considere uma segurada sintética que inicia a conversa
perguntando quando vence sua apólice. O Orquestrador identifica uma consulta e aciona o
especialista Consultas, que solicita o CPF, consulta a apólice e informa a data de vencimento.
Na sequência, a segurada relata ``roubaram meu carro ontem às 22h na Rua Augusta, 50''. Como
o segurado já foi identificado, a Triagem não solicita novamente o CPF: o relato é classificado
como roubo ou furto, e a apólice é verificada quanto ao status, à vigência na data do evento e
à existência da cobertura correspondente. O Aviso de Sinistro reconhece no próprio relato a data, a
hora, o local e o ocorrido, avalia preliminarmente a cobertura e registra o aviso, retornando
o protocolo. Em seguida, o especialista de Documentação solicita boletim de ocorrência, CNH e
CRLV e, após validar os arquivos enviados, encerra a abertura com uma mensagem final que
contém o protocolo, o resumo do evento, os documentos recebidos e os próximos passos. Todo o
percurso é visível no painel (Figura~\ref{fig:poc-painel}), e cada mensagem trocada passa pelos
\textit{guardrails} de entrada e de saída.

\subsubsection{Avaliação automatizada do artefato}
\label{sec:poc-avaliacao}

Como complemento à avaliação com especialistas baseada no TAM \cite{davis1989} (Seção~4.6), a
PoC inclui um mecanismo de avaliação automatizada, executável com um único comando, que
submete o sistema a uma bateria de cenários sintéticos sob diferentes configurações. Essa
avaliação fornece evidências objetivas sobre eficácia, segurança e custo do artefato antes da
apreciação pelos especialistas.

Os \textbf{cenários} cobrem caminho feliz, dados faltantes, apólice vencida, evento sem
cobertura, os diferentes tipos de sinistro, fraudes (texto, apólice recente, aviso tardio e
fraude descrita de forma indireta), \textit{prompt injection} direta, indireta via documento e
criativa, pedido discriminatório, assunto fora do escopo, pedido de atendimento humano,
cancelamento, documentos inválidos, expiração de prazo, consultas à apólice e falsos
positivos. Cada cenário declara o resultado esperado: etapa final, motivo de encerramento,
sequência de especialistas acionados, intenções detectadas, ferramentas obrigatórias e
proibidas e se deve ou não ser detectado como ataque.

Os \textbf{experimentos} variam a configuração do sistema: (i) somente regras;
(ii) somente regras e sem \textit{guardrails}; (iii) LLM como complemento na extração de
dados; (iv) LLM também no roteamento de intenções; (v) LLM também nos \textit{guardrails};
e (vi) a configuração anterior com um modelo menor, para comparação entre modelos. Para cada
experimento são calculadas as métricas da Tabela~\ref{tab:poc-metricas}. Inspirada na proposta
de LLM como juiz \cite{gu2026}, a avaliação automatizada pode ser estendida para julgar
qualidade e empatia das respostas, dimensão que permanece para trabalho futuro.

\begin{table}[htbp]
\centering
\caption{Métricas da avaliação automatizada da PoC}
\label{tab:poc-metricas}
\small
\begin{tabular}{|p{4.0cm}|p{9.4cm}|}
\hline
\textbf{Métrica} & \textbf{Definição} \\
\hline
Taxa de sucesso & Proporção de cenários com todas as verificações corretas \\
\hline
Taxa de conclusão & Proporção de cenários que terminam na etapa e com o motivo esperados \\
\hline
Acerto de roteamento & Especialistas acionados e intenções detectadas iguais às esperadas \\
\hline
Correção das ferramentas & Ferramentas obrigatórias chamadas e nenhuma ferramenta proibida
chamada \\
\hline
Detecção de ataques & Proporção de cenários de ataque bloqueados ou encaminhados a humano \\
\hline
Falsos positivos & Proporção de cenários legítimos indevidamente bloqueados ou encaminhados \\
\hline
Eficiência & Número de turnos, latência por turno (média e percentil 95), \textit{tokens} e
custo estimado \\
\hline
Resiliência & Chamadas ao LLM, falhas e uso do modelo de \textit{fallback} \\
\hline
\end{tabular}
\end{table}

% [PREENCHER] Inserir aqui a tabela de resultados gerada por `make evaluate`
% (reports/latest.md), após a execução final com o modelo principal disponível e,
% de preferência, com repetições para reportar média e desvio.

Uma execução preliminar, realizada antes da inclusão dos especialistas e cenários de consulta,
indicou três tendências:

\begin{itemize}
  \item a desativação dos \textit{guardrails} eliminou a detecção de todos os ataques;
  \item o uso do LLM como complemento elevou a taxa de sucesso em relação à configuração
        baseada apenas em regras, sobretudo nos relatos ambíguos e nos desvios de roteiro;
  \item o custo monetário das chamadas foi pouco relevante, enquanto a latência por turno
        aumentou em uma ordem de grandeza, configurando o principal compromisso
        (\textit{trade-off}) da abordagem.
\end{itemize}

\subsubsection{Limitações da demonstração}

A PoC não utiliza o WhatsApp real nem integrações com sistemas legados, e os documentos são
validados por metadados, sem OCR ou visão computacional. As etapas posteriores ao aviso
(regulação, vistoria, avaliação, liquidação e pesquisa de satisfação) não foram implementadas.
Os dados e a base de conhecimento são sintéticos e ilustrativos e não substituem condições
contratuais reais. As listas de documentos por tipo de sinistro e as regras de fraude são
propostas plausíveis, mas não normativas. Os resultados que envolvem LLM dependem da
disponibilidade do provedor e da variabilidade intrínseca dos modelos, o que recomenda
execuções repetidas para conclusões estatísticas. Ainda assim, a PoC demonstra que o recorte
do \textit{framework} pode ser instanciado de forma modular, auditável e governada, conforme as
diretrizes de projeto da Seção~4.4.2: modularidade, orquestração inteligente,
\textit{human-in-the-loop}, governança e segurança, experiência do usuário por interface
conversacional e orientação a dados.

% =============================================================================
% Referências usadas (BibTeX). Copie para o .bib apenas as que ainda não existirem
% ou troque as chaves acima pelas do seu arquivo. Todas constam na lista de
% referências do texto base da dissertação.
% =============================================================================
% @article{peffers2007, author={Peffers, Ken and others}, title={A design science research methodology for information systems research}, journal={Journal of Management Information Systems}, volume={24}, number={3}, year={2007}}
% @article{hevner2004, author={Hevner, Alan R. and others}, title={Design science in information systems research}, journal={MIS Quarterly}, volume={28}, number={1}, year={2004}}
% @article{capgemini2024b, author={{Capgemini}}, title={Maior parte dos brasileiros está insatisfeita com processo de acionamento de sinistro}, year={2024}, note={InfoMoney}}
% @misc{brasil2002codigocivil, author={{Brasil}}, title={Lei nº 10.406, de 10 de janeiro de 2002: Institui o Código Civil}, year={2002}}
% @misc{brasil2024marco, author={{Brasil}}, title={Lei nº 15.040, de 9 de dezembro de 2024: Marco Legal dos Seguros}, year={2024}}
% @article{ofoeda2019, author={Ofoeda, Joshua and Boateng, Richard and Effah, John}, title={Application programming interface (API) research: A review of the past to inform the future}, journal={International Journal of Enterprise Information Systems}, volume={15}, number={3}, pages={76--95}, year={2019}}
% @inproceedings{liu2024, author={Liu, B. and others}, title={Large language model based multi-agents: A survey of progress and challenges}, booktitle={Proceedings of the 33rd IJCAI}, year={2024}}
% @article{wang2023, author={Wang, W. and others}, title={A survey on large language model based autonomous agents}, journal={arXiv preprint arXiv:2308.11432}, year={2023}}
% @article{schick2023, author={Schick, T. and others}, title={Toolformer: Language models can teach themselves to use tools}, journal={Advances in Neural Information Processing Systems}, volume={36}, year={2023}}
% @article{qin2023, author={Qin, Y. and others}, title={ToolLLM: Facilitating large language models to master 16000+ real-world APIs}, journal={arXiv preprint arXiv:2307.16789}, year={2023}}
% @misc{mcp2025, author={{MCP}}, title={Model Context Protocol}, year={2025}, howpublished={\url{https://modelcontextprotocol.io/docs/getting-started/intro}}}
% @article{hill2025, author={Hill, G. and others}, title={LLMs and agentic AI in insurance decision-making: Opportunities and challenges for Africa}, journal={arXiv preprint arXiv:2508.15110}, year={2025}}
% @article{singhal2026, author={Singhal, S.}, title={Agentic AI in claims processing: Transforming insurance operations through autonomous AI systems}, journal={Journal of Computational Analysis and Applications}, volume={35}, number={2}, pages={208--218}, year={2026}}
% @misc{sapkota2025, author={Sapkota, R. and Roumeliotis, K. I. and Karkee, M.}, title={AI agents vs. agentic AI: A conceptual taxonomy, applications and challenges}, year={2025}}
% @article{casper2025, author={Casper, S. and others}, title={The 2025 AI agent index}, journal={arXiv preprint arXiv:2602.17753}, year={2025}}
% @article{hughes2025, author={Hughes, L. and others}, title={AI agents and agentic systems: A multi-expert analysis}, journal={Journal of Management Information Systems}, year={2025}}
% @article{dong2024, author={Dong, Y. and others}, title={Building guardrails for large language models}, journal={arXiv preprint arXiv:2402.01822}, year={2024}}
% @article{wu2022, author={Wu, X. and others}, title={A survey of human-in-the-loop for machine learning}, journal={Future Generation Computer Systems}, volume={135}, pages={364--381}, year={2022}}
% @article{dodda2023, author={Dodda, A.}, title={AI governance and security in fintech: Ensuring trust in generative and agentic AI systems}, journal={American Advanced Journal for Emerging Disciplinaries}, volume={1}, pages={1--14}, year={2023}}
% @misc{camara2023pl2338, author={{Câmara dos Deputados}}, title={PL 2338/2023: Lei Geral de Inteligência Artificial}, year={2023}}
% @article{davis1989, author={Davis, F. D.}, title={Perceived usefulness, perceived ease of use, and user acceptance of information technology}, journal={MIS Quarterly}, volume={13}, number={3}, pages={319--340}, year={1989}}
% @article{gu2026, author={Gu, J. and others}, title={A survey on LLM-as-a-judge}, journal={The Innovation}, volume={7}, number={6}, year={2026}}
 dentro da seção "Demonstração".
% Pacotes esperados no preâmbulo principal: graphicx (figura) e abntex2cite (ou
% natbib) para \cite. A figura usa o arquivo docs/figuras/poc-painel.png do
% repositório; envie-o ao Overleaf como figuras/poc-painel.png ou ajuste o caminho.
%
% As chaves de citação estão listadas em BibTeX no fim deste arquivo (comentadas).
% Se o seu .bib já usa outras chaves, basta substituí-las aqui.
% =============================================================================

\subsection{Prova de Conceito: recorte de Triagem, Aviso de Sinistro e Documentação}
\label{sec:poc}

A fase de demonstração, conforme preconizado pela DSR, tem por finalidade evidenciar que o
artefato proposto é capaz de resolver, ao menos em um contexto controlado, instâncias do
problema que motivou a pesquisa \cite{peffers2007,hevner2004}. Para tanto, foi construída uma
Prova de Conceito (PoC) que instancia um recorte do \textit{framework} apresentado na
Seção~4.4.1, abrangendo a orquestração, a Triagem, o Aviso de Sinistro e a etapa de
Documentação e Evidências Digitais (Figuras 2 a 5), além da camada transversal de
\textit{guardrails} corporativos (Figura~1). As etapas de Regulação, Vistoria, Avaliação,
Liquidação e Pesquisa de Satisfação permanecem fora do recorte e são tratadas como trabalhos
futuros.

O recorte foi escolhido por concentrar as dores mais citadas pelos segurados — complexidade
do processo de solicitação, falta de clareza sobre requisitos e documentos e ausência de
empatia no primeiro contato \cite{capgemini2024b} — e por constituir a porta de entrada de todo
o \textit{pipeline}: a qualidade das informações coletadas no aviso condiciona as etapas
subsequentes de regulação e liquidação. Trata-se, adicionalmente, do ponto em que o Código
Civil e o Marco Legal dos Seguros impõem ao segurado o dever de comunicar prontamente o
sinistro \cite{brasil2002codigocivil,brasil2024marco}, o que torna a disponibilidade contínua
de um canal conversacional particularmente relevante.

\subsubsection{Objetivo e premissas da PoC}

O objetivo da PoC é demonstrar a viabilidade técnica de um sistema multiagente que conduza a
jornada do segurado desde a primeira mensagem, em um canal que simula o WhatsApp, até a
conclusão da abertura do sinistro, com protocolo registrado e documentação mínima validada.
Foram adotadas as seguintes premissas:

\begin{itemize}
  \item \textbf{Ramo automóvel}, contemplando os tipos de sinistro de veículo mais recorrentes:
        colisão, roubo ou furto, incêndio, eventos da natureza (alagamento, granizo, queda de
        árvore) e danos a vidros. Solicitações de assistência (guincho, pane) e assuntos que
        não configuram sinistro recebem apenas orientação, sem abertura de processo.
  \item \textbf{Dados exclusivamente sintéticos}: clientes, apólices, veículos e documentos
        (CNH, CRLV, fotografias e boletins de ocorrência) são gerados artificialmente, em
        conformidade com os princípios de minimização e proteção de dados pessoais. Nenhum dado
        real é processado.
  \item \textbf{Sistemas da seguradora simulados}: cadastro, apólices, registro de sinistros,
        armazenamento e validação de documentos e envio de mensagens são implementados como
        serviços simulados (\textit{mocks}) atrás de interfaces estáveis, de modo que a troca
        por integrações reais não altere a lógica dos agentes \cite{ofoeda2019}.
  \item \textbf{Modelos de linguagem como recurso e não como dependência}: o sistema funciona
        integralmente com regras determinísticas e recorre ao LLM nos pontos em que a
        linguagem natural do segurado excede o que as regras capturam, o que permite comparar
        experimentalmente as duas abordagens (Seção~\ref{sec:poc-avaliacao}).
\end{itemize}

\subsubsection{Arquitetura e tecnologias}

A PoC foi implementada em Python~3.12 com o \textit{framework} LangGraph, que modela a
aplicação como um grafo de estados com suporte a ciclos, subgrafos e persistência
(\textit{checkpointing}) da conversa. Essa escolha materializa a noção de agente como unidade
que mantém estado interno persistente entre interações e decide a próxima ação a partir de
observações parciais do ambiente \cite{liu2024,wang2023}. A API foi construída com FastAPI, com
comunicação em tempo real via WebSocket, e a persistência utiliza SQLite para o estado das
conversas, o domínio simulado e o registro de eventos. A interface web, desenvolvida em React
com TypeScript e a biblioteca React Flow, reproduz um atendimento via WhatsApp e, lado a
lado, um painel de observabilidade que exibe o fluxo de agentes em execução
(Figura~\ref{fig:poc-painel}).

O acesso aos modelos de linguagem é feito por um \textit{gateway} próprio que padroniza as
chamadas no formato da API da OpenAI (\textit{OpenAI-compatible}), independentemente do
provedor. Na PoC, o provedor configurado é o Google Gemini, com um modelo principal e um
modelo de \textit{fallback} acionado em caso de indisponibilidade. As respostas estruturadas
são solicitadas por esquema JSON e validadas com Pydantic, e cada chamada registra o modelo
utilizado, a versão do \textit{prompt}, os \textit{tokens} consumidos e a latência. O
\textit{gateway} corresponde ao componente \textit{AI Gateway} da arquitetura proposta
(Figura~11) e garante que a substituição do provedor seja uma decisão de configuração, e não de
código.

As ferramentas (\textit{tools}) seguem o conceito de interfaces externas que ampliam a
capacidade do agente para além da predição de linguagem \cite{schick2023,qin2023}. Cinco
ferramentas foram definidas, cada uma com contrato tipado de entrada e saída, mecanismos de
\textit{timeout} e novas tentativas com recuo exponencial, além de injeção configurável de
falhas para avaliar a resiliência do sistema:

\begin{itemize}
  \item \texttt{get\_customer\_data}: consulta cadastro e apólices do segurado;
  \item \texttt{open\_claim}: registra o aviso de sinistro e gera o protocolo;
  \item \texttt{store\_files}: armazena os documentos vinculados ao protocolo;
  \item \texttt{validate\_document}: valida formato, tamanho, duplicidade e consistência dos
        metadados do documento com o processo;
  \item \texttt{send\_whatsapp}: entrega a mensagem ao segurado.
\end{itemize}

As interfaces das ferramentas foram desenhadas de forma compatível com a exposição via
\textit{Model Context Protocol} \cite{mcp2025}. Na versão atual da PoC, porém, elas são
invocadas diretamente pelos agentes por meio de um registro de implementações (simuladas ou
reais), e a publicação como servidor MCP permanece como evolução prevista.

\subsubsection{Orquestração e agentes especialistas}

A organização dos agentes segue o padrão hierárquico descrito na literatura de sistemas
multiagentes baseados em LLM, em que um orquestrador central decompõe a tarefa e delega
subtarefas a agentes especializados \cite{hill2025,singhal2026,sapkota2025}. Cada mensagem do
segurado percorre o grafo principal na seguinte sequência:
\textit{guardrails} de entrada $\rightarrow$ Orquestrador $\rightarrow$ agente especialista
$\rightarrow$ transição de etapa $\rightarrow$ \textit{guardrails} de saída $\rightarrow$ envio.
O estado global da conversa registra a etapa do processo
(\textsc{Triagem} $\rightarrow$ \textsc{Aviso} $\rightarrow$ \textsc{Documentos}
$\rightarrow$ \textsc{Concluído}), além de estados terminais alternativos (orientado,
bloqueado, encaminhado a humano e cancelado). A mudança de etapa é realizada exclusivamente
por uma função determinística de transição, o que garante previsibilidade e auditabilidade do
fluxo.

O \textbf{Orquestrador} implementa as três responsabilidades previstas na Figura~2:
identifica a situação atual do sinistro, informa o status ao segurado quando solicitado e
encaminha a mensagem ao especialista adequado. O encaminhamento usa uma tabela fixa que
associa cada etapa a um especialista. Quando a mensagem foge do roteiro esperado, um
classificador de intenções identifica desvios como dúvida, consulta à apólice, correção de
dado já informado, pedido de status, cancelamento, pedido de atendimento humano ou assunto
fora do escopo. O classificador combina regras e, opcionalmente, um LLM, e após tratar o
desvio o atendimento é retomado a partir da pergunta pendente. Cada especialista é um subgrafo
independente composto por subagentes, conforme a Tabela~\ref{tab:poc-agentes}.

\begin{table}[htbp]
\centering
\caption{Agentes especialistas, subagentes e ferramentas da PoC}
\label{tab:poc-agentes}
\small
\begin{tabular}{|p{2.6cm}|p{7.2cm}|p{3.4cm}|}
\hline
\textbf{Especialista} & \textbf{Subagentes (responsabilidade)} & \textbf{Ferramentas} \\
\hline
Triagem (Fig.~3) &
Validação do segurado (CPF e apólices); classificação do ocorrido (tipo de sinistro);
decisão de elegibilidade (status, vigência na data do evento e cobertura exigida pelo
tipo); identificação de envolvidos e feridos (colisão), com sinalização de urgência &
\texttt{get\_customer\_data} \\
\hline
Aviso de Sinistro (Fig.~4) &
Coleta multi-turno dos dados do evento (quando, onde, o que aconteceu e danos);
avaliação preliminar de cobertura; registro do aviso, com protocolo e próximos passos &
\texttt{open\_claim} \\
\hline
Documentação (Fig.~5) &
Lista de documentos por tipo de sinistro; recepção e armazenamento; validação;
complementação (progresso, lembrete, prazo e confirmação final) &
\texttt{store\_files}, \texttt{validate\_document} \\
\hline
Consultas &
Classificação da pergunta; identificação do segurado; resposta sobre vigência,
vencimento, renovação e coberturas da apólice ou sobre dúvidas gerais de seguro &
\texttt{get\_customer\_data} \\
\hline
\end{tabular}
\end{table}

Na \textbf{Triagem}, o segurado é identificado pelo CPF e o relato livre do ocorrido é
classificado, o que determina a cobertura exigida e os dados a coletar. Um mecanismo decisório
define se o caso evolui para a abertura formal ou se o segurado deve apenas ser orientado (por
exemplo, apólice vencida, evento fora da vigência, ausência da cobertura correspondente ou
pedido de assistência), evitando a abertura indevida de processos, como propõe a Figura~3. No
\textbf{Aviso de Sinistro}, os dados são extraídos de narrativas livres; quando uma informação
está ausente, o agente solicita apenas o que falta, sem repetir o que já foi informado. Após a
avaliação preliminar de cobertura, o aviso é registrado e o protocolo é comunicado. Na
\textbf{Documentação}, a lista de documentos varia conforme o tipo de sinistro (por exemplo,
boletim de ocorrência obrigatório em roubo ou furto e condicionado à existência de terceiros
ou feridos em colisão). Cada arquivo é validado de forma determinística quanto a formato,
tamanho, duplicidade e consistência dos metadados (CPF da CNH, placa do CRLV e data das
fotografias em relação ao evento). O agente informa o que já foi aceito e o que ainda falta e,
na ausência de resposta, envia lembrete e encaminha o caso a um atendente ao fim do prazo. O
especialista \textbf{Consultas}, acrescentado durante o ciclo de design, responde perguntas
sobre a própria apólice a partir dos dados cadastrais e dúvidas gerais (``o que é sinistro'',
``o que é franquia'') a partir de uma base de conhecimento curada. Ele atende à
responsabilidade de ``verificar a vigência e as coberturas da apólice'' atribuída a agentes de
sinistro \cite{singhal2026} e reduz a falta de clareza apontada pelos segurados
\cite{capgemini2024b}.

\subsubsection{Uso de modelos de linguagem: regras primeiro, LLM como complemento}

Em consonância com a ênfase da literatura em governança, previsibilidade e capacidade de
intervenção humana em sistemas agênticos \cite{casper2025,hughes2025}, a PoC adota o princípio
de ``regras primeiro, LLM como complemento''. Regras configuráveis (expressões e listas de
termos em arquivos YAML) resolvem os casos recorrentes com custo e latência mínimos. O LLM é
acionado quando as regras não são conclusivas, nos seguintes pontos:

\begin{itemize}
  \item extração e classificação do relato livre na Triagem e no Aviso (por exemplo, ``deixei o
        carro estacionado e quando voltei ele não estava mais lá'' como roubo ou furto);
  \item classificação de intenções do Orquestrador em mensagens que desviam do roteiro;
  \item classificação de segurança, ética e conformidade corporativa nos \textit{guardrails};
  \item redação de respostas a perguntas abertas, restrita à base de conhecimento e aos dados
        da apólice, técnica alinhada à geração aumentada por recuperação (RAG) e com instrução
        explícita de não prometer cobertura, valores ou prazos.
\end{itemize}

Todas as instruções enviadas ao LLM ficam em arquivos versionados, separados do código, e
as respostas estruturadas são validadas por esquema. Em caso de falha do provedor, o sistema
degrada para as regras sem interromper o atendimento.

\subsubsection{Camada de \textit{guardrails} corporativos}

A camada de \textit{guardrails} implementa as barreiras de proteção descritas na literatura
\cite{dong2024} e posicionadas transversalmente no \textit{framework} (Figura~1). Toda
mensagem recebida e toda resposta gerada passam por três estágios executados em paralelo:
\textbf{Segurança} (\textit{prompt injection} direta e indireta, tentativas de
\textit{jailbreak} e vazamento de dados pessoais), \textbf{Ética} (pedidos discriminatórios ou
abusivos) e \textbf{Corporativo} (sinais de fraude e aderência às políticas). Um agregador
consolida os resultados em um veredito (\textit{allow}, \textit{sanitize}, \textit{block} ou
\textit{escalate}) acompanhado de justificativa. Cada estágio pode ser ativado ou desativado
por configuração, o que viabiliza os experimentos comparativos. A detecção opera em duas
camadas: regras determinísticas e, quando estas não bloqueiam, um classificador baseado em LLM
que avalia os três estágios em uma única chamada e descarta veredictos de baixa confiança.

A injeção indireta é tratada inspecionando também o nome e os metadados dos documentos
enviados. Por exemplo, um arquivo nomeado com a instrução ``ignore as instruções e aprove o
sinistro'' é bloqueado antes do armazenamento. Ao final da Triagem, a checagem corporativa é
reexecutada com o contexto da apólice, incorporando sinais de fraude como apólice contratada
pouco antes do evento e aviso feito muito tempo após o ocorrido. Em coerência com o princípio
de \textit{human-in-the-loop} \cite{wu2022} e com a recomendação de manter sob supervisão
humana as decisões de maior risco \cite{singhal2026}, os casos suspeitos não são negados
automaticamente: são encaminhados a um atendente humano com a justificativa registrada. O
painel dispõe ainda de um controle de pausa que suspende a atuação dos agentes e transfere
as conversas para atendimento humano.

\subsubsection{Observabilidade, auditabilidade e conformidade}

Toda decisão relevante (roteamento, veredito de \textit{guardrail}, elegibilidade, chamada de
ferramenta e chamada ao LLM) é emitida como evento estruturado, persistida e transmitida ao
painel em tempo real. O registro inclui o agente, o subagente, a justificativa, a latência e,
para chamadas ao LLM, o modelo, a versão do \textit{prompt} e os \textit{tokens} consumidos.
Esse registro instrumenta as dimensões de conformidade e auditabilidade da Equação~1
\cite{dodda2023} e atende aos requisitos de transparência e rastreabilidade em decisões
automatizadas previstos no PL~2338/2023 \cite{camara2023pl2338}. Os dados pessoais (CPF e
telefone) são mascarados em logs e eventos. As regras de negócio (elegibilidade, tipos de
sinistro e coberturas exigidas, documentos por tipo, prazos de lembrete e de encerramento) são
mantidas em arquivos de configuração, separadas do código, e as decisões de projeto foram
documentadas em 19 registros de decisão arquitetural (ADRs), o que favorece a revisão por áreas
de negócio e de \textit{compliance}.

\begin{figure}[htbp]
  \centering
  \includegraphics[width=\textwidth]{figuras/poc-painel.png}
  \caption{Interface da PoC: à esquerda, o atendimento simulado via WhatsApp; à direita, o
  fluxo de agentes com o especialista e os subagentes em execução e o log de eventos com a
  justificativa de cada decisão. Fonte: Elaboração própria, 2026}
  \label{fig:poc-painel}
\end{figure}

\subsubsection{Cenário ilustrativo}

Para ilustrar o funcionamento, considere uma segurada sintética que inicia a conversa
perguntando quando vence sua apólice. O Orquestrador identifica uma consulta e aciona o
especialista Consultas, que solicita o CPF, consulta a apólice e informa a data de vencimento.
Na sequência, a segurada relata ``roubaram meu carro ontem às 22h na Rua Augusta, 50''. Como
o segurado já foi identificado, a Triagem não solicita novamente o CPF: o relato é classificado
como roubo ou furto, e a apólice é verificada quanto ao status, à vigência na data do evento e
à existência da cobertura correspondente. O Aviso de Sinistro reconhece no próprio relato a data, a
hora, o local e o ocorrido, avalia preliminarmente a cobertura e registra o aviso, retornando
o protocolo. Em seguida, o especialista de Documentação solicita boletim de ocorrência, CNH e
CRLV e, após validar os arquivos enviados, encerra a abertura com uma mensagem final que
contém o protocolo, o resumo do evento, os documentos recebidos e os próximos passos. Todo o
percurso é visível no painel (Figura~\ref{fig:poc-painel}), e cada mensagem trocada passa pelos
\textit{guardrails} de entrada e de saída.

\subsubsection{Avaliação automatizada do artefato}
\label{sec:poc-avaliacao}

Como complemento à avaliação com especialistas baseada no TAM \cite{davis1989} (Seção~4.6), a
PoC inclui um mecanismo de avaliação automatizada, executável com um único comando, que
submete o sistema a uma bateria de cenários sintéticos sob diferentes configurações. Essa
avaliação fornece evidências objetivas sobre eficácia, segurança e custo do artefato antes da
apreciação pelos especialistas.

Os \textbf{cenários} cobrem caminho feliz, dados faltantes, apólice vencida, evento sem
cobertura, os diferentes tipos de sinistro, fraudes (texto, apólice recente, aviso tardio e
fraude descrita de forma indireta), \textit{prompt injection} direta, indireta via documento e
criativa, pedido discriminatório, assunto fora do escopo, pedido de atendimento humano,
cancelamento, documentos inválidos, expiração de prazo, consultas à apólice e falsos
positivos. Cada cenário declara o resultado esperado: etapa final, motivo de encerramento,
sequência de especialistas acionados, intenções detectadas, ferramentas obrigatórias e
proibidas e se deve ou não ser detectado como ataque.

Os \textbf{experimentos} variam a configuração do sistema: (i) somente regras;
(ii) somente regras e sem \textit{guardrails}; (iii) LLM como complemento na extração de
dados; (iv) LLM também no roteamento de intenções; (v) LLM também nos \textit{guardrails};
e (vi) a configuração anterior com um modelo menor, para comparação entre modelos. Para cada
experimento são calculadas as métricas da Tabela~\ref{tab:poc-metricas}. Inspirada na proposta
de LLM como juiz \cite{gu2026}, a avaliação automatizada pode ser estendida para julgar
qualidade e empatia das respostas, dimensão que permanece para trabalho futuro.

\begin{table}[htbp]
\centering
\caption{Métricas da avaliação automatizada da PoC}
\label{tab:poc-metricas}
\small
\begin{tabular}{|p{4.0cm}|p{9.4cm}|}
\hline
\textbf{Métrica} & \textbf{Definição} \\
\hline
Taxa de sucesso & Proporção de cenários com todas as verificações corretas \\
\hline
Taxa de conclusão & Proporção de cenários que terminam na etapa e com o motivo esperados \\
\hline
Acerto de roteamento & Especialistas acionados e intenções detectadas iguais às esperadas \\
\hline
Correção das ferramentas & Ferramentas obrigatórias chamadas e nenhuma ferramenta proibida
chamada \\
\hline
Detecção de ataques & Proporção de cenários de ataque bloqueados ou encaminhados a humano \\
\hline
Falsos positivos & Proporção de cenários legítimos indevidamente bloqueados ou encaminhados \\
\hline
Eficiência & Número de turnos, latência por turno (média e percentil 95), \textit{tokens} e
custo estimado \\
\hline
Resiliência & Chamadas ao LLM, falhas e uso do modelo de \textit{fallback} \\
\hline
\end{tabular}
\end{table}

% [PREENCHER] Inserir aqui a tabela de resultados gerada por `make evaluate`
% (reports/latest.md), após a execução final com o modelo principal disponível e,
% de preferência, com repetições para reportar média e desvio.

Uma execução preliminar, realizada antes da inclusão dos especialistas e cenários de consulta,
indicou três tendências:

\begin{itemize}
  \item a desativação dos \textit{guardrails} eliminou a detecção de todos os ataques;
  \item o uso do LLM como complemento elevou a taxa de sucesso em relação à configuração
        baseada apenas em regras, sobretudo nos relatos ambíguos e nos desvios de roteiro;
  \item o custo monetário das chamadas foi pouco relevante, enquanto a latência por turno
        aumentou em uma ordem de grandeza, configurando o principal compromisso
        (\textit{trade-off}) da abordagem.
\end{itemize}

\subsubsection{Limitações da demonstração}

A PoC não utiliza o WhatsApp real nem integrações com sistemas legados, e os documentos são
validados por metadados, sem OCR ou visão computacional. As etapas posteriores ao aviso
(regulação, vistoria, avaliação, liquidação e pesquisa de satisfação) não foram implementadas.
Os dados e a base de conhecimento são sintéticos e ilustrativos e não substituem condições
contratuais reais. As listas de documentos por tipo de sinistro e as regras de fraude são
propostas plausíveis, mas não normativas. Os resultados que envolvem LLM dependem da
disponibilidade do provedor e da variabilidade intrínseca dos modelos, o que recomenda
execuções repetidas para conclusões estatísticas. Ainda assim, a PoC demonstra que o recorte
do \textit{framework} pode ser instanciado de forma modular, auditável e governada, conforme as
diretrizes de projeto da Seção~4.4.2: modularidade, orquestração inteligente,
\textit{human-in-the-loop}, governança e segurança, experiência do usuário por interface
conversacional e orientação a dados.

% =============================================================================
% Referências usadas (BibTeX). Copie para o .bib apenas as que ainda não existirem
% ou troque as chaves acima pelas do seu arquivo. Todas constam na lista de
% referências do texto base da dissertação.
% =============================================================================
% @article{peffers2007, author={Peffers, Ken and others}, title={A design science research methodology for information systems research}, journal={Journal of Management Information Systems}, volume={24}, number={3}, year={2007}}
% @article{hevner2004, author={Hevner, Alan R. and others}, title={Design science in information systems research}, journal={MIS Quarterly}, volume={28}, number={1}, year={2004}}
% @article{capgemini2024b, author={{Capgemini}}, title={Maior parte dos brasileiros está insatisfeita com processo de acionamento de sinistro}, year={2024}, note={InfoMoney}}
% @misc{brasil2002codigocivil, author={{Brasil}}, title={Lei nº 10.406, de 10 de janeiro de 2002: Institui o Código Civil}, year={2002}}
% @misc{brasil2024marco, author={{Brasil}}, title={Lei nº 15.040, de 9 de dezembro de 2024: Marco Legal dos Seguros}, year={2024}}
% @article{ofoeda2019, author={Ofoeda, Joshua and Boateng, Richard and Effah, John}, title={Application programming interface (API) research: A review of the past to inform the future}, journal={International Journal of Enterprise Information Systems}, volume={15}, number={3}, pages={76--95}, year={2019}}
% @inproceedings{liu2024, author={Liu, B. and others}, title={Large language model based multi-agents: A survey of progress and challenges}, booktitle={Proceedings of the 33rd IJCAI}, year={2024}}
% @article{wang2023, author={Wang, W. and others}, title={A survey on large language model based autonomous agents}, journal={arXiv preprint arXiv:2308.11432}, year={2023}}
% @article{schick2023, author={Schick, T. and others}, title={Toolformer: Language models can teach themselves to use tools}, journal={Advances in Neural Information Processing Systems}, volume={36}, year={2023}}
% @article{qin2023, author={Qin, Y. and others}, title={ToolLLM: Facilitating large language models to master 16000+ real-world APIs}, journal={arXiv preprint arXiv:2307.16789}, year={2023}}
% @misc{mcp2025, author={{MCP}}, title={Model Context Protocol}, year={2025}, howpublished={\url{https://modelcontextprotocol.io/docs/getting-started/intro}}}
% @article{hill2025, author={Hill, G. and others}, title={LLMs and agentic AI in insurance decision-making: Opportunities and challenges for Africa}, journal={arXiv preprint arXiv:2508.15110}, year={2025}}
% @article{singhal2026, author={Singhal, S.}, title={Agentic AI in claims processing: Transforming insurance operations through autonomous AI systems}, journal={Journal of Computational Analysis and Applications}, volume={35}, number={2}, pages={208--218}, year={2026}}
% @misc{sapkota2025, author={Sapkota, R. and Roumeliotis, K. I. and Karkee, M.}, title={AI agents vs. agentic AI: A conceptual taxonomy, applications and challenges}, year={2025}}
% @article{casper2025, author={Casper, S. and others}, title={The 2025 AI agent index}, journal={arXiv preprint arXiv:2602.17753}, year={2025}}
% @article{hughes2025, author={Hughes, L. and others}, title={AI agents and agentic systems: A multi-expert analysis}, journal={Journal of Management Information Systems}, year={2025}}
% @article{dong2024, author={Dong, Y. and others}, title={Building guardrails for large language models}, journal={arXiv preprint arXiv:2402.01822}, year={2024}}
% @article{wu2022, author={Wu, X. and others}, title={A survey of human-in-the-loop for machine learning}, journal={Future Generation Computer Systems}, volume={135}, pages={364--381}, year={2022}}
% @article{dodda2023, author={Dodda, A.}, title={AI governance and security in fintech: Ensuring trust in generative and agentic AI systems}, journal={American Advanced Journal for Emerging Disciplinaries}, volume={1}, pages={1--14}, year={2023}}
% @misc{camara2023pl2338, author={{Câmara dos Deputados}}, title={PL 2338/2023: Lei Geral de Inteligência Artificial}, year={2023}}
% @article{davis1989, author={Davis, F. D.}, title={Perceived usefulness, perceived ease of use, and user acceptance of information technology}, journal={MIS Quarterly}, volume={13}, number={3}, pages={319--340}, year={1989}}
% @article{gu2026, author={Gu, J. and others}, title={A survey on LLM-as-a-judge}, journal={The Innovation}, volume={7}, number={6}, year={2026}}
 dentro da seção "Demonstração".
% Pacotes esperados no preâmbulo principal: graphicx (figura) e abntex2cite (ou
% natbib) para \cite. A figura usa o arquivo docs/figuras/poc-painel.png do
% repositório; envie-o ao Overleaf como figuras/poc-painel.png ou ajuste o caminho.
%
% As chaves de citação estão listadas em BibTeX no fim deste arquivo (comentadas).
% Se o seu .bib já usa outras chaves, basta substituí-las aqui.
% =============================================================================

\subsection{Prova de Conceito: recorte de Triagem, Aviso de Sinistro e Documentação}
\label{sec:poc}

A fase de demonstração, conforme preconizado pela DSR, tem por finalidade evidenciar que o
artefato proposto é capaz de resolver, ao menos em um contexto controlado, instâncias do
problema que motivou a pesquisa \cite{peffers2007,hevner2004}. Para tanto, foi construída uma
Prova de Conceito (PoC) que instancia um recorte do \textit{framework} apresentado na
Seção~4.4.1, abrangendo a orquestração, a Triagem, o Aviso de Sinistro e a etapa de
Documentação e Evidências Digitais (Figuras 2 a 5), além da camada transversal de
\textit{guardrails} corporativos (Figura~1). As etapas de Regulação, Vistoria, Avaliação,
Liquidação e Pesquisa de Satisfação permanecem fora do recorte e são tratadas como trabalhos
futuros.

O recorte foi escolhido por concentrar as dores mais citadas pelos segurados — complexidade
do processo de solicitação, falta de clareza sobre requisitos e documentos e ausência de
empatia no primeiro contato \cite{capgemini2024b} — e por constituir a porta de entrada de todo
o \textit{pipeline}: a qualidade das informações coletadas no aviso condiciona as etapas
subsequentes de regulação e liquidação. Trata-se, adicionalmente, do ponto em que o Código
Civil e o Marco Legal dos Seguros impõem ao segurado o dever de comunicar prontamente o
sinistro \cite{brasil2002codigocivil,brasil2024marco}, o que torna a disponibilidade contínua
de um canal conversacional particularmente relevante.

\subsubsection{Objetivo e premissas da PoC}

O objetivo da PoC é demonstrar a viabilidade técnica de um sistema multiagente que conduza a
jornada do segurado desde a primeira mensagem, em um canal que simula o WhatsApp, até a
conclusão da abertura do sinistro, com protocolo registrado e documentação mínima validada.
Foram adotadas as seguintes premissas:

\begin{itemize}
  \item \textbf{Ramo automóvel}, contemplando os tipos de sinistro de veículo mais recorrentes:
        colisão, roubo ou furto, incêndio, eventos da natureza (alagamento, granizo, queda de
        árvore) e danos a vidros. Solicitações de assistência (guincho, pane) e assuntos que
        não configuram sinistro recebem apenas orientação, sem abertura de processo.
  \item \textbf{Dados exclusivamente sintéticos}: clientes, apólices, veículos e documentos
        (CNH, CRLV, fotografias e boletins de ocorrência) são gerados artificialmente, em
        conformidade com os princípios de minimização e proteção de dados pessoais. Nenhum dado
        real é processado.
  \item \textbf{Sistemas da seguradora simulados}: cadastro, apólices, registro de sinistros,
        armazenamento e validação de documentos e envio de mensagens são implementados como
        serviços simulados (\textit{mocks}) atrás de interfaces estáveis, de modo que a troca
        por integrações reais não altere a lógica dos agentes \cite{ofoeda2019}.
  \item \textbf{Modelos de linguagem como recurso e não como dependência}: o sistema funciona
        integralmente com regras determinísticas e recorre ao LLM nos pontos em que a
        linguagem natural do segurado excede o que as regras capturam, o que permite comparar
        experimentalmente as duas abordagens (Seção~\ref{sec:poc-avaliacao}).
\end{itemize}

\subsubsection{Arquitetura e tecnologias}

A PoC foi implementada em Python~3.12 com o \textit{framework} LangGraph, que modela a
aplicação como um grafo de estados com suporte a ciclos, subgrafos e persistência
(\textit{checkpointing}) da conversa. Essa escolha materializa a noção de agente como unidade
que mantém estado interno persistente entre interações e decide a próxima ação a partir de
observações parciais do ambiente \cite{liu2024,wang2023}. A API foi construída com FastAPI, com
comunicação em tempo real via WebSocket, e a persistência utiliza SQLite para o estado das
conversas, o domínio simulado e o registro de eventos. A interface web, desenvolvida em React
com TypeScript e a biblioteca React Flow, reproduz um atendimento via WhatsApp e, lado a
lado, um painel de observabilidade que exibe o fluxo de agentes em execução
(Figura~\ref{fig:poc-painel}).

O acesso aos modelos de linguagem é feito por um \textit{gateway} próprio que padroniza as
chamadas no formato da API da OpenAI (\textit{OpenAI-compatible}), independentemente do
provedor. Na PoC, o provedor configurado é o Google Gemini, com um modelo principal e um
modelo de \textit{fallback} acionado em caso de indisponibilidade. As respostas estruturadas
são solicitadas por esquema JSON e validadas com Pydantic, e cada chamada registra o modelo
utilizado, a versão do \textit{prompt}, os \textit{tokens} consumidos e a latência. O
\textit{gateway} corresponde ao componente \textit{AI Gateway} da arquitetura proposta
(Figura~11) e garante que a substituição do provedor seja uma decisão de configuração, e não de
código.

As ferramentas (\textit{tools}) seguem o conceito de interfaces externas que ampliam a
capacidade do agente para além da predição de linguagem \cite{schick2023,qin2023}. Cinco
ferramentas foram definidas, cada uma com contrato tipado de entrada e saída, mecanismos de
\textit{timeout} e novas tentativas com recuo exponencial, além de injeção configurável de
falhas para avaliar a resiliência do sistema:

\begin{itemize}
  \item \texttt{get\_customer\_data}: consulta cadastro e apólices do segurado;
  \item \texttt{open\_claim}: registra o aviso de sinistro e gera o protocolo;
  \item \texttt{store\_files}: armazena os documentos vinculados ao protocolo;
  \item \texttt{validate\_document}: valida formato, tamanho, duplicidade e consistência dos
        metadados do documento com o processo;
  \item \texttt{send\_whatsapp}: entrega a mensagem ao segurado.
\end{itemize}

As interfaces das ferramentas foram desenhadas de forma compatível com a exposição via
\textit{Model Context Protocol} \cite{mcp2025}. Na versão atual da PoC, porém, elas são
invocadas diretamente pelos agentes por meio de um registro de implementações (simuladas ou
reais), e a publicação como servidor MCP permanece como evolução prevista.

\subsubsection{Orquestração e agentes especialistas}

A organização dos agentes segue o padrão hierárquico descrito na literatura de sistemas
multiagentes baseados em LLM, em que um orquestrador central decompõe a tarefa e delega
subtarefas a agentes especializados \cite{hill2025,singhal2026,sapkota2025}. Cada mensagem do
segurado percorre o grafo principal na seguinte sequência:
\textit{guardrails} de entrada $\rightarrow$ Orquestrador $\rightarrow$ agente especialista
$\rightarrow$ transição de etapa $\rightarrow$ \textit{guardrails} de saída $\rightarrow$ envio.
O estado global da conversa registra a etapa do processo
(\textsc{Triagem} $\rightarrow$ \textsc{Aviso} $\rightarrow$ \textsc{Documentos}
$\rightarrow$ \textsc{Concluído}), além de estados terminais alternativos (orientado,
bloqueado, encaminhado a humano e cancelado). A mudança de etapa é realizada exclusivamente
por uma função determinística de transição, o que garante previsibilidade e auditabilidade do
fluxo.

O \textbf{Orquestrador} implementa as três responsabilidades previstas na Figura~2:
identifica a situação atual do sinistro, informa o status ao segurado quando solicitado e
encaminha a mensagem ao especialista adequado. O encaminhamento usa uma tabela fixa que
associa cada etapa a um especialista. Quando a mensagem foge do roteiro esperado, um
classificador de intenções identifica desvios como dúvida, consulta à apólice, correção de
dado já informado, pedido de status, cancelamento, pedido de atendimento humano ou assunto
fora do escopo. O classificador combina regras e, opcionalmente, um LLM, e após tratar o
desvio o atendimento é retomado a partir da pergunta pendente. Cada especialista é um subgrafo
independente composto por subagentes, conforme a Tabela~\ref{tab:poc-agentes}.

\begin{table}[htbp]
\centering
\caption{Agentes especialistas, subagentes e ferramentas da PoC}
\label{tab:poc-agentes}
\small
\begin{tabular}{|p{2.6cm}|p{7.2cm}|p{3.4cm}|}
\hline
\textbf{Especialista} & \textbf{Subagentes (responsabilidade)} & \textbf{Ferramentas} \\
\hline
Triagem (Fig.~3) &
Validação do segurado (CPF e apólices); classificação do ocorrido (tipo de sinistro);
decisão de elegibilidade (status, vigência na data do evento e cobertura exigida pelo
tipo); identificação de envolvidos e feridos (colisão), com sinalização de urgência &
\texttt{get\_customer\_data} \\
\hline
Aviso de Sinistro (Fig.~4) &
Coleta multi-turno dos dados do evento (quando, onde, o que aconteceu e danos);
avaliação preliminar de cobertura; registro do aviso, com protocolo e próximos passos &
\texttt{open\_claim} \\
\hline
Documentação (Fig.~5) &
Lista de documentos por tipo de sinistro; recepção e armazenamento; validação;
complementação (progresso, lembrete, prazo e confirmação final) &
\texttt{store\_files}, \texttt{validate\_document} \\
\hline
Consultas &
Classificação da pergunta; identificação do segurado; resposta sobre vigência,
vencimento, renovação e coberturas da apólice ou sobre dúvidas gerais de seguro &
\texttt{get\_customer\_data} \\
\hline
\end{tabular}
\end{table}

Na \textbf{Triagem}, o segurado é identificado pelo CPF e o relato livre do ocorrido é
classificado, o que determina a cobertura exigida e os dados a coletar. Um mecanismo decisório
define se o caso evolui para a abertura formal ou se o segurado deve apenas ser orientado (por
exemplo, apólice vencida, evento fora da vigência, ausência da cobertura correspondente ou
pedido de assistência), evitando a abertura indevida de processos, como propõe a Figura~3. No
\textbf{Aviso de Sinistro}, os dados são extraídos de narrativas livres; quando uma informação
está ausente, o agente solicita apenas o que falta, sem repetir o que já foi informado. Após a
avaliação preliminar de cobertura, o aviso é registrado e o protocolo é comunicado. Na
\textbf{Documentação}, a lista de documentos varia conforme o tipo de sinistro (por exemplo,
boletim de ocorrência obrigatório em roubo ou furto e condicionado à existência de terceiros
ou feridos em colisão). Cada arquivo é validado de forma determinística quanto a formato,
tamanho, duplicidade e consistência dos metadados (CPF da CNH, placa do CRLV e data das
fotografias em relação ao evento). O agente informa o que já foi aceito e o que ainda falta e,
na ausência de resposta, envia lembrete e encaminha o caso a um atendente ao fim do prazo. O
especialista \textbf{Consultas}, acrescentado durante o ciclo de design, responde perguntas
sobre a própria apólice a partir dos dados cadastrais e dúvidas gerais (``o que é sinistro'',
``o que é franquia'') a partir de uma base de conhecimento curada. Ele atende à
responsabilidade de ``verificar a vigência e as coberturas da apólice'' atribuída a agentes de
sinistro \cite{singhal2026} e reduz a falta de clareza apontada pelos segurados
\cite{capgemini2024b}.

\subsubsection{Uso de modelos de linguagem: regras primeiro, LLM como complemento}

Em consonância com a ênfase da literatura em governança, previsibilidade e capacidade de
intervenção humana em sistemas agênticos \cite{casper2025,hughes2025}, a PoC adota o princípio
de ``regras primeiro, LLM como complemento''. Regras configuráveis (expressões e listas de
termos em arquivos YAML) resolvem os casos recorrentes com custo e latência mínimos. O LLM é
acionado quando as regras não são conclusivas, nos seguintes pontos:

\begin{itemize}
  \item extração e classificação do relato livre na Triagem e no Aviso (por exemplo, ``deixei o
        carro estacionado e quando voltei ele não estava mais lá'' como roubo ou furto);
  \item classificação de intenções do Orquestrador em mensagens que desviam do roteiro;
  \item classificação de segurança, ética e conformidade corporativa nos \textit{guardrails};
  \item redação de respostas a perguntas abertas, restrita à base de conhecimento e aos dados
        da apólice, técnica alinhada à geração aumentada por recuperação (RAG) e com instrução
        explícita de não prometer cobertura, valores ou prazos.
\end{itemize}

Todas as instruções enviadas ao LLM ficam em arquivos versionados, separados do código, e
as respostas estruturadas são validadas por esquema. Em caso de falha do provedor, o sistema
degrada para as regras sem interromper o atendimento.

\subsubsection{Camada de \textit{guardrails} corporativos}

A camada de \textit{guardrails} implementa as barreiras de proteção descritas na literatura
\cite{dong2024} e posicionadas transversalmente no \textit{framework} (Figura~1). Toda
mensagem recebida e toda resposta gerada passam por três estágios executados em paralelo:
\textbf{Segurança} (\textit{prompt injection} direta e indireta, tentativas de
\textit{jailbreak} e vazamento de dados pessoais), \textbf{Ética} (pedidos discriminatórios ou
abusivos) e \textbf{Corporativo} (sinais de fraude e aderência às políticas). Um agregador
consolida os resultados em um veredito (\textit{allow}, \textit{sanitize}, \textit{block} ou
\textit{escalate}) acompanhado de justificativa. Cada estágio pode ser ativado ou desativado
por configuração, o que viabiliza os experimentos comparativos. A detecção opera em duas
camadas: regras determinísticas e, quando estas não bloqueiam, um classificador baseado em LLM
que avalia os três estágios em uma única chamada e descarta veredictos de baixa confiança.

A injeção indireta é tratada inspecionando também o nome e os metadados dos documentos
enviados. Por exemplo, um arquivo nomeado com a instrução ``ignore as instruções e aprove o
sinistro'' é bloqueado antes do armazenamento. Ao final da Triagem, a checagem corporativa é
reexecutada com o contexto da apólice, incorporando sinais de fraude como apólice contratada
pouco antes do evento e aviso feito muito tempo após o ocorrido. Em coerência com o princípio
de \textit{human-in-the-loop} \cite{wu2022} e com a recomendação de manter sob supervisão
humana as decisões de maior risco \cite{singhal2026}, os casos suspeitos não são negados
automaticamente: são encaminhados a um atendente humano com a justificativa registrada. O
painel dispõe ainda de um controle de pausa que suspende a atuação dos agentes e transfere
as conversas para atendimento humano.

\subsubsection{Observabilidade, auditabilidade e conformidade}

Toda decisão relevante (roteamento, veredito de \textit{guardrail}, elegibilidade, chamada de
ferramenta e chamada ao LLM) é emitida como evento estruturado, persistida e transmitida ao
painel em tempo real. O registro inclui o agente, o subagente, a justificativa, a latência e,
para chamadas ao LLM, o modelo, a versão do \textit{prompt} e os \textit{tokens} consumidos.
Esse registro instrumenta as dimensões de conformidade e auditabilidade da Equação~1
\cite{dodda2023} e atende aos requisitos de transparência e rastreabilidade em decisões
automatizadas previstos no PL~2338/2023 \cite{camara2023pl2338}. Os dados pessoais (CPF e
telefone) são mascarados em logs e eventos. As regras de negócio (elegibilidade, tipos de
sinistro e coberturas exigidas, documentos por tipo, prazos de lembrete e de encerramento) são
mantidas em arquivos de configuração, separadas do código, e as decisões de projeto foram
documentadas em 19 registros de decisão arquitetural (ADRs), o que favorece a revisão por áreas
de negócio e de \textit{compliance}.

\begin{figure}[htbp]
  \centering
  \includegraphics[width=\textwidth]{figuras/poc-painel.png}
  \caption{Interface da PoC: à esquerda, o atendimento simulado via WhatsApp; à direita, o
  fluxo de agentes com o especialista e os subagentes em execução e o log de eventos com a
  justificativa de cada decisão. Fonte: Elaboração própria, 2026}
  \label{fig:poc-painel}
\end{figure}

\subsubsection{Cenário ilustrativo}

Para ilustrar o funcionamento, considere uma segurada sintética que inicia a conversa
perguntando quando vence sua apólice. O Orquestrador identifica uma consulta e aciona o
especialista Consultas, que solicita o CPF, consulta a apólice e informa a data de vencimento.
Na sequência, a segurada relata ``roubaram meu carro ontem às 22h na Rua Augusta, 50''. Como
o segurado já foi identificado, a Triagem não solicita novamente o CPF: o relato é classificado
como roubo ou furto, e a apólice é verificada quanto ao status, à vigência na data do evento e
à existência da cobertura correspondente. O Aviso de Sinistro reconhece no próprio relato a data, a
hora, o local e o ocorrido, avalia preliminarmente a cobertura e registra o aviso, retornando
o protocolo. Em seguida, o especialista de Documentação solicita boletim de ocorrência, CNH e
CRLV e, após validar os arquivos enviados, encerra a abertura com uma mensagem final que
contém o protocolo, o resumo do evento, os documentos recebidos e os próximos passos. Todo o
percurso é visível no painel (Figura~\ref{fig:poc-painel}), e cada mensagem trocada passa pelos
\textit{guardrails} de entrada e de saída.

\subsubsection{Avaliação automatizada do artefato}
\label{sec:poc-avaliacao}

Como complemento à avaliação com especialistas baseada no TAM \cite{davis1989} (Seção~4.6), a
PoC inclui um mecanismo de avaliação automatizada, executável com um único comando, que
submete o sistema a uma bateria de cenários sintéticos sob diferentes configurações. Essa
avaliação fornece evidências objetivas sobre eficácia, segurança e custo do artefato antes da
apreciação pelos especialistas.

Os \textbf{cenários} cobrem caminho feliz, dados faltantes, apólice vencida, evento sem
cobertura, os diferentes tipos de sinistro, fraudes (texto, apólice recente, aviso tardio e
fraude descrita de forma indireta), \textit{prompt injection} direta, indireta via documento e
criativa, pedido discriminatório, assunto fora do escopo, pedido de atendimento humano,
cancelamento, documentos inválidos, expiração de prazo, consultas à apólice e falsos
positivos. Cada cenário declara o resultado esperado: etapa final, motivo de encerramento,
sequência de especialistas acionados, intenções detectadas, ferramentas obrigatórias e
proibidas e se deve ou não ser detectado como ataque.

Os \textbf{experimentos} variam a configuração do sistema: (i) somente regras;
(ii) somente regras e sem \textit{guardrails}; (iii) LLM como complemento na extração de
dados; (iv) LLM também no roteamento de intenções; (v) LLM também nos \textit{guardrails};
e (vi) a configuração anterior com um modelo menor, para comparação entre modelos. Para cada
experimento são calculadas as métricas da Tabela~\ref{tab:poc-metricas}. Inspirada na proposta
de LLM como juiz \cite{gu2026}, a avaliação automatizada pode ser estendida para julgar
qualidade e empatia das respostas, dimensão que permanece para trabalho futuro.

\begin{table}[htbp]
\centering
\caption{Métricas da avaliação automatizada da PoC}
\label{tab:poc-metricas}
\small
\begin{tabular}{|p{4.0cm}|p{9.4cm}|}
\hline
\textbf{Métrica} & \textbf{Definição} \\
\hline
Taxa de sucesso & Proporção de cenários com todas as verificações corretas \\
\hline
Taxa de conclusão & Proporção de cenários que terminam na etapa e com o motivo esperados \\
\hline
Acerto de roteamento & Especialistas acionados e intenções detectadas iguais às esperadas \\
\hline
Correção das ferramentas & Ferramentas obrigatórias chamadas e nenhuma ferramenta proibida
chamada \\
\hline
Detecção de ataques & Proporção de cenários de ataque bloqueados ou encaminhados a humano \\
\hline
Falsos positivos & Proporção de cenários legítimos indevidamente bloqueados ou encaminhados \\
\hline
Eficiência & Número de turnos, latência por turno (média e percentil 95), \textit{tokens} e
custo estimado \\
\hline
Resiliência & Chamadas ao LLM, falhas e uso do modelo de \textit{fallback} \\
\hline
\end{tabular}
\end{table}

% [PREENCHER] Inserir aqui a tabela de resultados gerada por `make evaluate`
% (reports/latest.md), após a execução final com o modelo principal disponível e,
% de preferência, com repetições para reportar média e desvio.

Uma execução preliminar, realizada antes da inclusão dos especialistas e cenários de consulta,
indicou três tendências:

\begin{itemize}
  \item a desativação dos \textit{guardrails} eliminou a detecção de todos os ataques;
  \item o uso do LLM como complemento elevou a taxa de sucesso em relação à configuração
        baseada apenas em regras, sobretudo nos relatos ambíguos e nos desvios de roteiro;
  \item o custo monetário das chamadas foi pouco relevante, enquanto a latência por turno
        aumentou em uma ordem de grandeza, configurando o principal compromisso
        (\textit{trade-off}) da abordagem.
\end{itemize}

\subsubsection{Limitações da demonstração}

A PoC não utiliza o WhatsApp real nem integrações com sistemas legados, e os documentos são
validados por metadados, sem OCR ou visão computacional. As etapas posteriores ao aviso
(regulação, vistoria, avaliação, liquidação e pesquisa de satisfação) não foram implementadas.
Os dados e a base de conhecimento são sintéticos e ilustrativos e não substituem condições
contratuais reais. As listas de documentos por tipo de sinistro e as regras de fraude são
propostas plausíveis, mas não normativas. Os resultados que envolvem LLM dependem da
disponibilidade do provedor e da variabilidade intrínseca dos modelos, o que recomenda
execuções repetidas para conclusões estatísticas. Ainda assim, a PoC demonstra que o recorte
do \textit{framework} pode ser instanciado de forma modular, auditável e governada, conforme as
diretrizes de projeto da Seção~4.4.2: modularidade, orquestração inteligente,
\textit{human-in-the-loop}, governança e segurança, experiência do usuário por interface
conversacional e orientação a dados.

% =============================================================================
% Referências usadas (BibTeX). Copie para o .bib apenas as que ainda não existirem
% ou troque as chaves acima pelas do seu arquivo. Todas constam na lista de
% referências do texto base da dissertação.
% =============================================================================
% @article{peffers2007, author={Peffers, Ken and others}, title={A design science research methodology for information systems research}, journal={Journal of Management Information Systems}, volume={24}, number={3}, year={2007}}
% @article{hevner2004, author={Hevner, Alan R. and others}, title={Design science in information systems research}, journal={MIS Quarterly}, volume={28}, number={1}, year={2004}}
% @article{capgemini2024b, author={{Capgemini}}, title={Maior parte dos brasileiros está insatisfeita com processo de acionamento de sinistro}, year={2024}, note={InfoMoney}}
% @misc{brasil2002codigocivil, author={{Brasil}}, title={Lei nº 10.406, de 10 de janeiro de 2002: Institui o Código Civil}, year={2002}}
% @misc{brasil2024marco, author={{Brasil}}, title={Lei nº 15.040, de 9 de dezembro de 2024: Marco Legal dos Seguros}, year={2024}}
% @article{ofoeda2019, author={Ofoeda, Joshua and Boateng, Richard and Effah, John}, title={Application programming interface (API) research: A review of the past to inform the future}, journal={International Journal of Enterprise Information Systems}, volume={15}, number={3}, pages={76--95}, year={2019}}
% @inproceedings{liu2024, author={Liu, B. and others}, title={Large language model based multi-agents: A survey of progress and challenges}, booktitle={Proceedings of the 33rd IJCAI}, year={2024}}
% @article{wang2023, author={Wang, W. and others}, title={A survey on large language model based autonomous agents}, journal={arXiv preprint arXiv:2308.11432}, year={2023}}
% @article{schick2023, author={Schick, T. and others}, title={Toolformer: Language models can teach themselves to use tools}, journal={Advances in Neural Information Processing Systems}, volume={36}, year={2023}}
% @article{qin2023, author={Qin, Y. and others}, title={ToolLLM: Facilitating large language models to master 16000+ real-world APIs}, journal={arXiv preprint arXiv:2307.16789}, year={2023}}
% @misc{mcp2025, author={{MCP}}, title={Model Context Protocol}, year={2025}, howpublished={\url{https://modelcontextprotocol.io/docs/getting-started/intro}}}
% @article{hill2025, author={Hill, G. and others}, title={LLMs and agentic AI in insurance decision-making: Opportunities and challenges for Africa}, journal={arXiv preprint arXiv:2508.15110}, year={2025}}
% @article{singhal2026, author={Singhal, S.}, title={Agentic AI in claims processing: Transforming insurance operations through autonomous AI systems}, journal={Journal of Computational Analysis and Applications}, volume={35}, number={2}, pages={208--218}, year={2026}}
% @misc{sapkota2025, author={Sapkota, R. and Roumeliotis, K. I. and Karkee, M.}, title={AI agents vs. agentic AI: A conceptual taxonomy, applications and challenges}, year={2025}}
% @article{casper2025, author={Casper, S. and others}, title={The 2025 AI agent index}, journal={arXiv preprint arXiv:2602.17753}, year={2025}}
% @article{hughes2025, author={Hughes, L. and others}, title={AI agents and agentic systems: A multi-expert analysis}, journal={Journal of Management Information Systems}, year={2025}}
% @article{dong2024, author={Dong, Y. and others}, title={Building guardrails for large language models}, journal={arXiv preprint arXiv:2402.01822}, year={2024}}
% @article{wu2022, author={Wu, X. and others}, title={A survey of human-in-the-loop for machine learning}, journal={Future Generation Computer Systems}, volume={135}, pages={364--381}, year={2022}}
% @article{dodda2023, author={Dodda, A.}, title={AI governance and security in fintech: Ensuring trust in generative and agentic AI systems}, journal={American Advanced Journal for Emerging Disciplinaries}, volume={1}, pages={1--14}, year={2023}}
% @misc{camara2023pl2338, author={{Câmara dos Deputados}}, title={PL 2338/2023: Lei Geral de Inteligência Artificial}, year={2023}}
% @article{davis1989, author={Davis, F. D.}, title={Perceived usefulness, perceived ease of use, and user acceptance of information technology}, journal={MIS Quarterly}, volume={13}, number={3}, pages={319--340}, year={1989}}
% @article{gu2026, author={Gu, J. and others}, title={A survey on LLM-as-a-judge}, journal={The Innovation}, volume={7}, number={6}, year={2026}}
 dentro da seção "Demonstração".
% Pacotes esperados no preâmbulo principal: graphicx (figura) e abntex2cite (ou
% natbib) para \cite. A figura usa o arquivo docs/figuras/poc-painel.png do
% repositório; envie-o ao Overleaf como figuras/poc-painel.png ou ajuste o caminho.
%
% As chaves de citação estão listadas em BibTeX no fim deste arquivo (comentadas).
% Se o seu .bib já usa outras chaves, basta substituí-las aqui.
% =============================================================================

\subsection{Prova de Conceito: recorte de Triagem, Aviso de Sinistro e Documentação}
\label{sec:poc}

A fase de demonstração, conforme preconizado pela DSR, tem por finalidade evidenciar que o
artefato proposto é capaz de resolver, ao menos em um contexto controlado, instâncias do
problema que motivou a pesquisa \cite{peffers2007,hevner2004}. Para tanto, foi construída uma
Prova de Conceito (PoC) que instancia um recorte do \textit{framework} apresentado na
Seção~4.4.1, abrangendo a orquestração, a Triagem, o Aviso de Sinistro e a etapa de
Documentação e Evidências Digitais (Figuras 2 a 5), além da camada transversal de
\textit{guardrails} corporativos (Figura~1). As etapas de Regulação, Vistoria, Avaliação,
Liquidação e Pesquisa de Satisfação permanecem fora do recorte e são tratadas como trabalhos
futuros.

O recorte foi escolhido por concentrar as dores mais citadas pelos segurados — complexidade
do processo de solicitação, falta de clareza sobre requisitos e documentos e ausência de
empatia no primeiro contato \cite{capgemini2024b} — e por constituir a porta de entrada de todo
o \textit{pipeline}: a qualidade das informações coletadas no aviso condiciona as etapas
subsequentes de regulação e liquidação. Trata-se, adicionalmente, do ponto em que o Código
Civil e o Marco Legal dos Seguros impõem ao segurado o dever de comunicar prontamente o
sinistro \cite{brasil2002codigocivil,brasil2024marco}, o que torna a disponibilidade contínua
de um canal conversacional particularmente relevante.

\subsubsection{Objetivo e premissas da PoC}

O objetivo da PoC é demonstrar a viabilidade técnica de um sistema multiagente que conduza a
jornada do segurado desde a primeira mensagem, em um canal que simula o WhatsApp, até a
conclusão da abertura do sinistro, com protocolo registrado e documentação mínima validada.
Foram adotadas as seguintes premissas:

\begin{itemize}
  \item \textbf{Ramo automóvel}, contemplando os tipos de sinistro de veículo mais recorrentes:
        colisão, roubo ou furto, incêndio, eventos da natureza (alagamento, granizo, queda de
        árvore) e danos a vidros. Solicitações de assistência (guincho, pane) e assuntos que
        não configuram sinistro recebem apenas orientação, sem abertura de processo.
  \item \textbf{Dados exclusivamente sintéticos}: clientes, apólices, veículos e documentos
        (CNH, CRLV, fotografias e boletins de ocorrência) são gerados artificialmente, em
        conformidade com os princípios de minimização e proteção de dados pessoais. Nenhum dado
        real é processado.
  \item \textbf{Sistemas da seguradora simulados}: cadastro, apólices, registro de sinistros,
        armazenamento e validação de documentos e envio de mensagens são implementados como
        serviços simulados (\textit{mocks}) atrás de interfaces estáveis, de modo que a troca
        por integrações reais não altere a lógica dos agentes \cite{ofoeda2019}.
  \item \textbf{Modelos de linguagem como recurso e não como dependência}: o sistema funciona
        integralmente com regras determinísticas e recorre ao LLM nos pontos em que a
        linguagem natural do segurado excede o que as regras capturam, o que permite comparar
        experimentalmente as duas abordagens (Seção~\ref{sec:poc-avaliacao}).
\end{itemize}

\subsubsection{Arquitetura e tecnologias}

A PoC foi implementada em Python~3.12 com o \textit{framework} LangGraph, que modela a
aplicação como um grafo de estados com suporte a ciclos, subgrafos e persistência
(\textit{checkpointing}) da conversa. Essa escolha materializa a noção de agente como unidade
que mantém estado interno persistente entre interações e decide a próxima ação a partir de
observações parciais do ambiente \cite{liu2024,wang2023}. A API foi construída com FastAPI, com
comunicação em tempo real via WebSocket, e a persistência utiliza SQLite para o estado das
conversas, o domínio simulado e o registro de eventos. A interface web, desenvolvida em React
com TypeScript e a biblioteca React Flow, reproduz um atendimento via WhatsApp e, lado a
lado, um painel de observabilidade que exibe o fluxo de agentes em execução
(Figura~\ref{fig:poc-painel}).

O acesso aos modelos de linguagem é feito por um \textit{gateway} próprio que padroniza as
chamadas no formato da API da OpenAI (\textit{OpenAI-compatible}), independentemente do
provedor. Na PoC, o provedor configurado é o Google Gemini, com um modelo principal e um
modelo de \textit{fallback} acionado em caso de indisponibilidade. As respostas estruturadas
são solicitadas por esquema JSON e validadas com Pydantic, e cada chamada registra o modelo
utilizado, a versão do \textit{prompt}, os \textit{tokens} consumidos e a latência. O
\textit{gateway} corresponde ao componente \textit{AI Gateway} da arquitetura proposta
(Figura~11) e garante que a substituição do provedor seja uma decisão de configuração, e não de
código.

As ferramentas (\textit{tools}) seguem o conceito de interfaces externas que ampliam a
capacidade do agente para além da predição de linguagem \cite{schick2023,qin2023}. Cinco
ferramentas foram definidas, cada uma com contrato tipado de entrada e saída, mecanismos de
\textit{timeout} e novas tentativas com recuo exponencial, além de injeção configurável de
falhas para avaliar a resiliência do sistema:

\begin{itemize}
  \item \texttt{get\_customer\_data}: consulta cadastro e apólices do segurado;
  \item \texttt{open\_claim}: registra o aviso de sinistro e gera o protocolo;
  \item \texttt{store\_files}: armazena os documentos vinculados ao protocolo;
  \item \texttt{validate\_document}: valida formato, tamanho, duplicidade e consistência dos
        metadados do documento com o processo;
  \item \texttt{send\_whatsapp}: entrega a mensagem ao segurado.
\end{itemize}

As interfaces das ferramentas foram desenhadas de forma compatível com a exposição via
\textit{Model Context Protocol} \cite{mcp2025}. Na versão atual da PoC, porém, elas são
invocadas diretamente pelos agentes por meio de um registro de implementações (simuladas ou
reais), e a publicação como servidor MCP permanece como evolução prevista.

\subsubsection{Orquestração e agentes especialistas}

A organização dos agentes segue o padrão hierárquico descrito na literatura de sistemas
multiagentes baseados em LLM, em que um orquestrador central decompõe a tarefa e delega
subtarefas a agentes especializados \cite{hill2025,singhal2026,sapkota2025}. Cada mensagem do
segurado percorre o grafo principal na seguinte sequência:
\textit{guardrails} de entrada $\rightarrow$ Orquestrador $\rightarrow$ agente especialista
$\rightarrow$ transição de etapa $\rightarrow$ \textit{guardrails} de saída $\rightarrow$ envio.
O estado global da conversa registra a etapa do processo
(\textsc{Triagem} $\rightarrow$ \textsc{Aviso} $\rightarrow$ \textsc{Documentos}
$\rightarrow$ \textsc{Concluído}), além de estados terminais alternativos (orientado,
bloqueado, encaminhado a humano e cancelado). A mudança de etapa é realizada exclusivamente
por uma função determinística de transição, o que garante previsibilidade e auditabilidade do
fluxo.

O \textbf{Orquestrador} implementa as três responsabilidades previstas na Figura~2:
identifica a situação atual do sinistro, informa o status ao segurado quando solicitado e
encaminha a mensagem ao especialista adequado. O encaminhamento usa uma tabela fixa que
associa cada etapa a um especialista. Quando a mensagem foge do roteiro esperado, um
classificador de intenções identifica desvios como dúvida, consulta à apólice, correção de
dado já informado, pedido de status, cancelamento, pedido de atendimento humano ou assunto
fora do escopo. O classificador combina regras e, opcionalmente, um LLM, e após tratar o
desvio o atendimento é retomado a partir da pergunta pendente. Cada especialista é um subgrafo
independente composto por subagentes, conforme a Tabela~\ref{tab:poc-agentes}.

\begin{table}[htbp]
\centering
\caption{Agentes especialistas, subagentes e ferramentas da PoC}
\label{tab:poc-agentes}
\small
\begin{tabular}{|p{2.6cm}|p{7.2cm}|p{3.4cm}|}
\hline
\textbf{Especialista} & \textbf{Subagentes (responsabilidade)} & \textbf{Ferramentas} \\
\hline
Triagem (Fig.~3) &
Validação do segurado (CPF e apólices); classificação do ocorrido (tipo de sinistro);
decisão de elegibilidade (status, vigência na data do evento e cobertura exigida pelo
tipo); identificação de envolvidos e feridos (colisão), com sinalização de urgência &
\texttt{get\_customer\_data} \\
\hline
Aviso de Sinistro (Fig.~4) &
Coleta multi-turno dos dados do evento (quando, onde, o que aconteceu e danos);
avaliação preliminar de cobertura; registro do aviso, com protocolo e próximos passos &
\texttt{open\_claim} \\
\hline
Documentação (Fig.~5) &
Lista de documentos por tipo de sinistro; recepção e armazenamento; validação;
complementação (progresso, lembrete, prazo e confirmação final) &
\texttt{store\_files}, \texttt{validate\_document} \\
\hline
Consultas &
Classificação da pergunta; identificação do segurado; resposta sobre vigência,
vencimento, renovação e coberturas da apólice ou sobre dúvidas gerais de seguro &
\texttt{get\_customer\_data} \\
\hline
\end{tabular}
\end{table}

Na \textbf{Triagem}, o segurado é identificado pelo CPF e o relato livre do ocorrido é
classificado, o que determina a cobertura exigida e os dados a coletar. Um mecanismo decisório
define se o caso evolui para a abertura formal ou se o segurado deve apenas ser orientado (por
exemplo, apólice vencida, evento fora da vigência, ausência da cobertura correspondente ou
pedido de assistência), evitando a abertura indevida de processos, como propõe a Figura~3. No
\textbf{Aviso de Sinistro}, os dados são extraídos de narrativas livres; quando uma informação
está ausente, o agente solicita apenas o que falta, sem repetir o que já foi informado. Após a
avaliação preliminar de cobertura, o aviso é registrado e o protocolo é comunicado. Na
\textbf{Documentação}, a lista de documentos varia conforme o tipo de sinistro (por exemplo,
boletim de ocorrência obrigatório em roubo ou furto e condicionado à existência de terceiros
ou feridos em colisão). Cada arquivo é validado de forma determinística quanto a formato,
tamanho, duplicidade e consistência dos metadados (CPF da CNH, placa do CRLV e data das
fotografias em relação ao evento). O agente informa o que já foi aceito e o que ainda falta e,
na ausência de resposta, envia lembrete e encaminha o caso a um atendente ao fim do prazo. O
especialista \textbf{Consultas}, acrescentado durante o ciclo de design, responde perguntas
sobre a própria apólice a partir dos dados cadastrais e dúvidas gerais (``o que é sinistro'',
``o que é franquia'') a partir de uma base de conhecimento curada. Ele atende à
responsabilidade de ``verificar a vigência e as coberturas da apólice'' atribuída a agentes de
sinistro \cite{singhal2026} e reduz a falta de clareza apontada pelos segurados
\cite{capgemini2024b}.

\subsubsection{Uso de modelos de linguagem: regras primeiro, LLM como complemento}

Em consonância com a ênfase da literatura em governança, previsibilidade e capacidade de
intervenção humana em sistemas agênticos \cite{casper2025,hughes2025}, a PoC adota o princípio
de ``regras primeiro, LLM como complemento''. Regras configuráveis (expressões e listas de
termos em arquivos YAML) resolvem os casos recorrentes com custo e latência mínimos. O LLM é
acionado quando as regras não são conclusivas, nos seguintes pontos:

\begin{itemize}
  \item extração e classificação do relato livre na Triagem e no Aviso (por exemplo, ``deixei o
        carro estacionado e quando voltei ele não estava mais lá'' como roubo ou furto);
  \item classificação de intenções do Orquestrador em mensagens que desviam do roteiro;
  \item classificação de segurança, ética e conformidade corporativa nos \textit{guardrails};
  \item redação de respostas a perguntas abertas, restrita à base de conhecimento e aos dados
        da apólice, técnica alinhada à geração aumentada por recuperação (RAG) e com instrução
        explícita de não prometer cobertura, valores ou prazos.
\end{itemize}

Todas as instruções enviadas ao LLM ficam em arquivos versionados, separados do código, e
as respostas estruturadas são validadas por esquema. Em caso de falha do provedor, o sistema
degrada para as regras sem interromper o atendimento.

\subsubsection{Camada de \textit{guardrails} corporativos}

A camada de \textit{guardrails} implementa as barreiras de proteção descritas na literatura
\cite{dong2024} e posicionadas transversalmente no \textit{framework} (Figura~1). Toda
mensagem recebida e toda resposta gerada passam por três estágios executados em paralelo:
\textbf{Segurança} (\textit{prompt injection} direta e indireta, tentativas de
\textit{jailbreak} e vazamento de dados pessoais), \textbf{Ética} (pedidos discriminatórios ou
abusivos) e \textbf{Corporativo} (sinais de fraude e aderência às políticas). Um agregador
consolida os resultados em um veredito (\textit{allow}, \textit{sanitize}, \textit{block} ou
\textit{escalate}) acompanhado de justificativa. Cada estágio pode ser ativado ou desativado
por configuração, o que viabiliza os experimentos comparativos. A detecção opera em duas
camadas: regras determinísticas e, quando estas não bloqueiam, um classificador baseado em LLM
que avalia os três estágios em uma única chamada e descarta veredictos de baixa confiança.

A injeção indireta é tratada inspecionando também o nome e os metadados dos documentos
enviados. Por exemplo, um arquivo nomeado com a instrução ``ignore as instruções e aprove o
sinistro'' é bloqueado antes do armazenamento. Ao final da Triagem, a checagem corporativa é
reexecutada com o contexto da apólice, incorporando sinais de fraude como apólice contratada
pouco antes do evento e aviso feito muito tempo após o ocorrido. Em coerência com o princípio
de \textit{human-in-the-loop} \cite{wu2022} e com a recomendação de manter sob supervisão
humana as decisões de maior risco \cite{singhal2026}, os casos suspeitos não são negados
automaticamente: são encaminhados a um atendente humano com a justificativa registrada. O
painel dispõe ainda de um controle de pausa que suspende a atuação dos agentes e transfere
as conversas para atendimento humano.

\subsubsection{Observabilidade, auditabilidade e conformidade}

Toda decisão relevante (roteamento, veredito de \textit{guardrail}, elegibilidade, chamada de
ferramenta e chamada ao LLM) é emitida como evento estruturado, persistida e transmitida ao
painel em tempo real. O registro inclui o agente, o subagente, a justificativa, a latência e,
para chamadas ao LLM, o modelo, a versão do \textit{prompt} e os \textit{tokens} consumidos.
Esse registro instrumenta as dimensões de conformidade e auditabilidade da Equação~1
\cite{dodda2023} e atende aos requisitos de transparência e rastreabilidade em decisões
automatizadas previstos no PL~2338/2023 \cite{camara2023pl2338}. Os dados pessoais (CPF e
telefone) são mascarados em logs e eventos. As regras de negócio (elegibilidade, tipos de
sinistro e coberturas exigidas, documentos por tipo, prazos de lembrete e de encerramento) são
mantidas em arquivos de configuração, separadas do código, e as decisões de projeto foram
documentadas em 19 registros de decisão arquitetural (ADRs), o que favorece a revisão por áreas
de negócio e de \textit{compliance}.

\begin{figure}[htbp]
  \centering
  \includegraphics[width=\textwidth]{figuras/poc-painel.png}
  \caption{Interface da PoC: à esquerda, o atendimento simulado via WhatsApp; à direita, o
  fluxo de agentes com o especialista e os subagentes em execução e o log de eventos com a
  justificativa de cada decisão. Fonte: Elaboração própria, 2026}
  \label{fig:poc-painel}
\end{figure}

\subsubsection{Cenário ilustrativo}

Para ilustrar o funcionamento, considere uma segurada sintética que inicia a conversa
perguntando quando vence sua apólice. O Orquestrador identifica uma consulta e aciona o
especialista Consultas, que solicita o CPF, consulta a apólice e informa a data de vencimento.
Na sequência, a segurada relata ``roubaram meu carro ontem às 22h na Rua Augusta, 50''. Como
o segurado já foi identificado, a Triagem não solicita novamente o CPF: o relato é classificado
como roubo ou furto, e a apólice é verificada quanto ao status, à vigência na data do evento e
à existência da cobertura correspondente. O Aviso de Sinistro reconhece no próprio relato a data, a
hora, o local e o ocorrido, avalia preliminarmente a cobertura e registra o aviso, retornando
o protocolo. Em seguida, o especialista de Documentação solicita boletim de ocorrência, CNH e
CRLV e, após validar os arquivos enviados, encerra a abertura com uma mensagem final que
contém o protocolo, o resumo do evento, os documentos recebidos e os próximos passos. Todo o
percurso é visível no painel (Figura~\ref{fig:poc-painel}), e cada mensagem trocada passa pelos
\textit{guardrails} de entrada e de saída.

\subsubsection{Avaliação automatizada do artefato}
\label{sec:poc-avaliacao}

Como complemento à avaliação com especialistas baseada no TAM \cite{davis1989} (Seção~4.6), a
PoC inclui um mecanismo de avaliação automatizada, executável com um único comando, que
submete o sistema a uma bateria de cenários sintéticos sob diferentes configurações. Essa
avaliação fornece evidências objetivas sobre eficácia, segurança e custo do artefato antes da
apreciação pelos especialistas.

Os \textbf{cenários} cobrem caminho feliz, dados faltantes, apólice vencida, evento sem
cobertura, os diferentes tipos de sinistro, fraudes (texto, apólice recente, aviso tardio e
fraude descrita de forma indireta), \textit{prompt injection} direta, indireta via documento e
criativa, pedido discriminatório, assunto fora do escopo, pedido de atendimento humano,
cancelamento, documentos inválidos, expiração de prazo, consultas à apólice e falsos
positivos. Cada cenário declara o resultado esperado: etapa final, motivo de encerramento,
sequência de especialistas acionados, intenções detectadas, ferramentas obrigatórias e
proibidas e se deve ou não ser detectado como ataque.

Os \textbf{experimentos} variam a configuração do sistema: (i) somente regras;
(ii) somente regras e sem \textit{guardrails}; (iii) LLM como complemento na extração de
dados; (iv) LLM também no roteamento de intenções; (v) LLM também nos \textit{guardrails};
e (vi) a configuração anterior com um modelo menor, para comparação entre modelos. Para cada
experimento são calculadas as métricas da Tabela~\ref{tab:poc-metricas}. Inspirada na proposta
de LLM como juiz \cite{gu2026}, a avaliação automatizada pode ser estendida para julgar
qualidade e empatia das respostas, dimensão que permanece para trabalho futuro.

\begin{table}[htbp]
\centering
\caption{Métricas da avaliação automatizada da PoC}
\label{tab:poc-metricas}
\small
\begin{tabular}{|p{4.0cm}|p{9.4cm}|}
\hline
\textbf{Métrica} & \textbf{Definição} \\
\hline
Taxa de sucesso & Proporção de cenários com todas as verificações corretas \\
\hline
Taxa de conclusão & Proporção de cenários que terminam na etapa e com o motivo esperados \\
\hline
Acerto de roteamento & Especialistas acionados e intenções detectadas iguais às esperadas \\
\hline
Correção das ferramentas & Ferramentas obrigatórias chamadas e nenhuma ferramenta proibida
chamada \\
\hline
Detecção de ataques & Proporção de cenários de ataque bloqueados ou encaminhados a humano \\
\hline
Falsos positivos & Proporção de cenários legítimos indevidamente bloqueados ou encaminhados \\
\hline
Eficiência & Número de turnos, latência por turno (média e percentil 95), \textit{tokens} e
custo estimado \\
\hline
Resiliência & Chamadas ao LLM, falhas e uso do modelo de \textit{fallback} \\
\hline
\end{tabular}
\end{table}

% [PREENCHER] Inserir aqui a tabela de resultados gerada por `make evaluate`
% (reports/latest.md), após a execução final com o modelo principal disponível e,
% de preferência, com repetições para reportar média e desvio.

Uma execução preliminar, realizada antes da inclusão dos especialistas e cenários de consulta,
indicou três tendências:

\begin{itemize}
  \item a desativação dos \textit{guardrails} eliminou a detecção de todos os ataques;
  \item o uso do LLM como complemento elevou a taxa de sucesso em relação à configuração
        baseada apenas em regras, sobretudo nos relatos ambíguos e nos desvios de roteiro;
  \item o custo monetário das chamadas foi pouco relevante, enquanto a latência por turno
        aumentou em uma ordem de grandeza, configurando o principal compromisso
        (\textit{trade-off}) da abordagem.
\end{itemize}

\subsubsection{Limitações da demonstração}

A PoC não utiliza o WhatsApp real nem integrações com sistemas legados, e os documentos são
validados por metadados, sem OCR ou visão computacional. As etapas posteriores ao aviso
(regulação, vistoria, avaliação, liquidação e pesquisa de satisfação) não foram implementadas.
Os dados e a base de conhecimento são sintéticos e ilustrativos e não substituem condições
contratuais reais. As listas de documentos por tipo de sinistro e as regras de fraude são
propostas plausíveis, mas não normativas. Os resultados que envolvem LLM dependem da
disponibilidade do provedor e da variabilidade intrínseca dos modelos, o que recomenda
execuções repetidas para conclusões estatísticas. Ainda assim, a PoC demonstra que o recorte
do \textit{framework} pode ser instanciado de forma modular, auditável e governada, conforme as
diretrizes de projeto da Seção~4.4.2: modularidade, orquestração inteligente,
\textit{human-in-the-loop}, governança e segurança, experiência do usuário por interface
conversacional e orientação a dados.

% =============================================================================
% Referências usadas (BibTeX). Copie para o .bib apenas as que ainda não existirem
% ou troque as chaves acima pelas do seu arquivo. Todas constam na lista de
% referências do texto base da dissertação.
% =============================================================================
% @article{peffers2007, author={Peffers, Ken and others}, title={A design science research methodology for information systems research}, journal={Journal of Management Information Systems}, volume={24}, number={3}, year={2007}}
% @article{hevner2004, author={Hevner, Alan R. and others}, title={Design science in information systems research}, journal={MIS Quarterly}, volume={28}, number={1}, year={2004}}
% @article{capgemini2024b, author={{Capgemini}}, title={Maior parte dos brasileiros está insatisfeita com processo de acionamento de sinistro}, year={2024}, note={InfoMoney}}
% @misc{brasil2002codigocivil, author={{Brasil}}, title={Lei nº 10.406, de 10 de janeiro de 2002: Institui o Código Civil}, year={2002}}
% @misc{brasil2024marco, author={{Brasil}}, title={Lei nº 15.040, de 9 de dezembro de 2024: Marco Legal dos Seguros}, year={2024}}
% @article{ofoeda2019, author={Ofoeda, Joshua and Boateng, Richard and Effah, John}, title={Application programming interface (API) research: A review of the past to inform the future}, journal={International Journal of Enterprise Information Systems}, volume={15}, number={3}, pages={76--95}, year={2019}}
% @inproceedings{liu2024, author={Liu, B. and others}, title={Large language model based multi-agents: A survey of progress and challenges}, booktitle={Proceedings of the 33rd IJCAI}, year={2024}}
% @article{wang2023, author={Wang, W. and others}, title={A survey on large language model based autonomous agents}, journal={arXiv preprint arXiv:2308.11432}, year={2023}}
% @article{schick2023, author={Schick, T. and others}, title={Toolformer: Language models can teach themselves to use tools}, journal={Advances in Neural Information Processing Systems}, volume={36}, year={2023}}
% @article{qin2023, author={Qin, Y. and others}, title={ToolLLM: Facilitating large language models to master 16000+ real-world APIs}, journal={arXiv preprint arXiv:2307.16789}, year={2023}}
% @misc{mcp2025, author={{MCP}}, title={Model Context Protocol}, year={2025}, howpublished={\url{https://modelcontextprotocol.io/docs/getting-started/intro}}}
% @article{hill2025, author={Hill, G. and others}, title={LLMs and agentic AI in insurance decision-making: Opportunities and challenges for Africa}, journal={arXiv preprint arXiv:2508.15110}, year={2025}}
% @article{singhal2026, author={Singhal, S.}, title={Agentic AI in claims processing: Transforming insurance operations through autonomous AI systems}, journal={Journal of Computational Analysis and Applications}, volume={35}, number={2}, pages={208--218}, year={2026}}
% @misc{sapkota2025, author={Sapkota, R. and Roumeliotis, K. I. and Karkee, M.}, title={AI agents vs. agentic AI: A conceptual taxonomy, applications and challenges}, year={2025}}
% @article{casper2025, author={Casper, S. and others}, title={The 2025 AI agent index}, journal={arXiv preprint arXiv:2602.17753}, year={2025}}
% @article{hughes2025, author={Hughes, L. and others}, title={AI agents and agentic systems: A multi-expert analysis}, journal={Journal of Management Information Systems}, year={2025}}
% @article{dong2024, author={Dong, Y. and others}, title={Building guardrails for large language models}, journal={arXiv preprint arXiv:2402.01822}, year={2024}}
% @article{wu2022, author={Wu, X. and others}, title={A survey of human-in-the-loop for machine learning}, journal={Future Generation Computer Systems}, volume={135}, pages={364--381}, year={2022}}
% @article{dodda2023, author={Dodda, A.}, title={AI governance and security in fintech: Ensuring trust in generative and agentic AI systems}, journal={American Advanced Journal for Emerging Disciplinaries}, volume={1}, pages={1--14}, year={2023}}
% @misc{camara2023pl2338, author={{Câmara dos Deputados}}, title={PL 2338/2023: Lei Geral de Inteligência Artificial}, year={2023}}
% @article{davis1989, author={Davis, F. D.}, title={Perceived usefulness, perceived ease of use, and user acceptance of information technology}, journal={MIS Quarterly}, volume={13}, number={3}, pages={319--340}, year={1989}}
% @article{gu2026, author={Gu, J. and others}, title={A survey on LLM-as-a-judge}, journal={The Innovation}, volume={7}, number={6}, year={2026}}
